\documentclass[a4paper,11pt,oneside]{article}
\usepackage[T1]{fontenc}
\usepackage{newunicodechar}
\usepackage[utf8]{inputenc}
\usepackage[table]{xcolor}
\usepackage{url}
\usepackage{booktabs}
\usepackage{longtable}
\usepackage{caption}
\usepackage[sc,osf]{mathpazo}
\linespread{1.05}
\usepackage[authoryear]{natbib}
\usepackage{siunitx}
\usepackage{comment}
\usepackage{rotating}
\usepackage{epsfig}
\usepackage{adjustbox}
\usepackage{algorithmic}
\usepackage{amsmath}
\usepackage{threeparttable}
\usepackage{amstext}
\usepackage{amssymb}
\usepackage{shadow}

\usepackage{amsbsy}

\usepackage{bm}

\usepackage{multirow}

\usepackage{hanging}


\usepackage{soul}

\usepackage{sectsty}

\usepackage{amsmath}
%\usepackage{chicago}
\usepackage{geometry}
\usepackage{setspace}

\usepackage{array}
\usepackage{tabularx}
\usepackage{color}
\usepackage{xcolor}
\usepackage{eurosym}
\usepackage{ntheorem}
\usepackage{booktabs}

\usepackage{subcaption}
\usepackage{dcolumn}

\usepackage{appendix}

\usepackage{lineno}

\usepackage{enumerate}
\usepackage{lscape}
\usepackage{longtable}

\theoremseparator{:}
\newtheorem{hyp}{Hypothesis}
\usepackage{ulem}

\newcolumntype{L}{>{$}l<{$}} %
\newcolumntype{C}{>{$}c<{$}} %
\newcolumntype{R}{>{$}r<{$}} %
\usepackage{graphicx}
\usepackage{needspace}

\usepackage{float}

\usepackage{booktabs}
\usepackage{threeparttable}
\usepackage{siunitx}
\usepackage{adjustbox}


\usepackage{geometry}

% \usepackage{endnotes}
% \let\footnote=\endnote

\geometry{margin=2.5cm}

%Configuración siunitx: SIN redondeo global
\sisetup{
  group-separator = {,},
  group-minimum-digits = 4,
  round-mode = none           % Use 'none' instead of 'off'
}

\textwidth 162mm \textheight 233mm \topmargin -1.3cm
\oddsidemargin -0,1cm \pagestyle{plain}
\newcommand{\rojo}{\textcolor[rgb]{1.00,0.00,0.00}}


\newcommand{\azul}{\textcolor[rgb]{0.00,0.00,1.00}}


\newcommand{\verde}{\textcolor[rgb]{0.00,0.50,0.70}}


\newcommand{\marron}{\textcolor[rgb]{0.80,0.50,0.00}}



% color de Luisa:
\definecolor{luisa}{rgb}{0.97,0.17,0.93}
\newcommand{\Luisa}{\color{luisa}}

% color de GUI:
\definecolor{miguelangel}{rgb}{0.97,0.60,0.13}
\newcommand{\Miguelangel}{\color{miguelangel}}

% color de Emili:
\definecolor{emili}{rgb}{0.77,0.47,0.23}
\newcommand{\Emili}{\color{emili}}


\newcommand{\m}{Malmquist }

\newcommand{\mpi}{Malmquist productivity index }
\usepackage{pdflscape}

\newcommand{\mpis}{Malmquist productivity indices }

\newcommand{\dos}{\renewcommand{\baselinestretch}{2}}
\renewcommand{\baselinestretch}{1.35}
\newcommand{\uno}{\renewcommand{\baselinestretch}{1}}

\renewcommand{\captionlabelfont}{\normalsize \bfseries}
\renewcommand{\captionfont}{\normalfont}

\newcommand\solidrule[1][1cm]{\rule[0.5ex]{#1}{.1pt}}
\newcommand\dashedrule{\mbox{%
  \solidrule[1mm]\hspace{.75mm}\solidrule[1mm]\hspace{.75mm}\solidrule[1mm]\hspace{.75mm}\solidrule[1mm]\hspace{.75mm}\solidrule[1mm]\hspace{.75mm}\solidrule[1mm]}}

\makeatletter
\def\@seccntformat#1{\csname the#1\endcsname.\quad}
\makeatother


\sectionfont{\large\bf}
\subsectionfont{\normalsize\bf}


\theorembodyfont{\sl}
\usepackage{rotating}

\newtheorem{theorem}{Theorem}
\newtheorem{acknowledgement}{Acknowledgement}
\newtheorem{algorithm}{Algorithm}
\newtheorem{axiom}{Axiom}
\newtheorem{case}{Case}
\newtheorem{claim}{Claim}
\newtheorem{conclusion}{Conclusion}
\newtheorem{condition}{Condition}
\newtheorem{conjecture}{Conjecture}
\newtheorem{corollary}{Corollary}
\newtheorem{criterion}{Criterion}
\newtheorem{definition}{Definition}

\newtheorem{example}{Example}
\newtheorem{exercise}{Exercise}
\newtheorem{lemma}{Lemma}
\newtheorem{notation}{Notation}
\newtheorem{problem}{Problem}
\newtheorem{proposition}{Proposition}
%\newtheorem{remark}{Remark}
\newtheorem{remark}{Nota}
\newtheorem{solution}{Solution}
\newtheorem{summary}{Summary}

\newenvironment{proof}[1][Proof]{\noindent\textbf{#1.} }{\ \rule{0.5em}{0.5em}}



\renewcommand{\baselinestretch}{1.35}

\hyphenation{ma-kers}
%\documentclass[a4paper,10pt]{article}
%\usepackage{amsmath}

\begin{document}

% \title{\textbf{Clustering Left Behind Places in Spain: A Multidimensional Municipality Approach}}

% Rethinking Left-Behindness: A Multidimensional Territorial Approach to Spanish Municipalities
% Y PROPUESTA 2

\title{\textbf{Spatial Context and the Typology of Left-Behind Places: Toward Systemic Place-Based Policies?}}

\author{\normalsize  Luisa Alam\'a-Sabater \\ \textsl{\footnotesize Universitat Jaume I and IIDL} \and  \normalsize Miguel \'A. M\'arquez \\ \textsl{\footnotesize Universidad de Extremadura}  \and \normalsize Guillem Rodilla\\ \textsl{\footnotesize Universitat Politècnica de València}
  \and     \normalsize Emili Tortosa-Ausina \\ \textsl{\footnotesize Universitat
Jaume I, IIDL and Ivie}}
% \date{February 2025}



\bibliographystyle{apalike}


\vspace{\stretch{1}}


\begin{spacing}{1}


\maketitle

\thispagestyle{empty}




\begin{abstract}

{\color{red} (Sólo un párrafo)}


This paper aims to analyse how demographic, socioeconomic, and second-nature geographic factors jointly shape the emergence, spatial distribution, and heterogeneity of left-behind places (LBPs) in Spain, while drawing implications for the design of more systemic and place-sensitive regional development policies. The proposed methodological framework aims to explain both municipalities’ proximity to vulnerability thresholds and their probability of belonging to specific LBP categories.

To this end, it develops an empirical strategy using a set of demographic, socioeconomic and spatial indicators at the municipal level. The results reveal substantial heterogeneity in left-behind dynamics both across and within clustering analysis, uncovering multiple forms of municipal vulnerability. Building on this classification, the paper combines a frontier analysis and a multinomial logit model. The frontier analysis measures the distance of each municipality to cluster-defining thresholds, capturing how close localities are to transitioning into more vulnerable states. In parallel, the multinomial model estimates the probability of belonging to each LBP category relative to non-left-behind municipalities, identifying the key drivers of these trajectories.


A central contribution of the paper is to highlight the role of second-nature geography in shaping these processes. Variables capturing accessibility to urban and semiurban centres and market potential emerge as critical factors, not only in structuring the clusters themselves but also in explaining municipalities’ position relative to vulnerability frontiers and their probability of belonging to specific LBP types. This underscores the importance of agglomeration forces and spatial interactions in the emergence and persistence of territorial inequalities. From a policy perspective, the findings support a shift towards what we call systemic place-based policies, which seek to reshape spatial structures by altering the second-nature geographic preconditions for economic activity.
\end{abstract}



\vspace{0.6cm}

\noindent {\textbf{Key words and phrases}: cluster analysis, depopulation, left-behind places, rural, urban}
\\

\noindent {\textbf{JEL Classifications}: C3, O18, O21, R1, R23, R3}
\\




\begin{sloppypar}
\noindent {\textbf{Communications to}: Emili Tortosa-Ausina, Departament d'Economia, Universitat Jaume I, Campus del Riu Sec, 12071 Castell\'o de la Plana, Spain. Tel.: +34 964387168, e-mail:} \texttt{tortosa@uji.es}.
\end{sloppypar}






\end{spacing}

\pagebreak




\section{Introduction}
\label{sec:intro}


%{\color{red} EMILI: motivaci\'on ser\'a cluster analysis to identify and characterise LBP given the richness of the municipal information in Spain (in terms of number of municipalities, and information available). Is such detailed information available in other regions? Also LBP/depopulation in Spain is important.}


\Emili

\begin{sloppypar}





Over the past decades, European regional policy \citep{fratesi2023regional} and the Cohesion Policy \citep{rodriguez2021does,di_caro_fratesi2022one,bourdin2025impacts} have aimed to reduce territorial disparities and promote balanced development across regions. Nevertheless, many scholars have pointed out that, despite decades of significant financial transfers, Cohesion Policy has often struggled to reverse the long-term trend of divergence between core European regions and those on the periphery, leading to growing calls for fundamental reforms \citep{rodriguez_pose2018revenge}. The ongoing policy debate \citep{rodriguez2025forging} reinforces these tensions: while the policy maintains its core mission, it faces the risk of fragmentation as it is increasingly tasked with addressing new strategic priorities like defence and competitiveness \citep{ECA2025}. Addressing this policy gap requires moving beyond macro-level diagnosis to understand the localized and multidimensional mechanisms of exclusion that define the specific nature of these development traps. Moreover, this situation points to the limitations of conventional place-based policies \citep{barca2009agenda,fratesi2023regional} and motivates the consideration of what we term ``systemic place-based policies.'' Such policies do not merely respond to existing spatial structures but actively seek to reshape them by altering the second-nature geographic preconditions\footnote{Second-nature geography refers to the set of spatial conditions shaped by economic forces and institutions (such as infrastructure, accessibility, and agglomeration economies) that influence the distribution of economic activity beyond natural endowments \citep[see, among others,][]{fujita2004new}. This paper contributes to the literature by analysing how the interaction between economic forces and spatial factors, as captured by second-nature geography, affects left-behindness. Hereafter, geography is conceptualized as ‘second-nature’ geography unless stated otherwise.} for economic activity. Systemic place-based policies can be defined as forward-looking, place-based capacity-building policies that contextualize local structural conditions while deliberately investing in enabling infrastructures which may exceed current demographic or economic needs, but are considered necessary to unlock future development trajectories in lagging territories.



\end{sloppypar}

{\color{red} No sé si estos dos párrafos conectan bien; pasar de cohesion a left-behind hay que hacerlo mejor. El párrafo de arriba también es muy largo, podría ser la mitad, el lector se pierde un poco. ¿Qué pasa si lo quitamos y empezamos directamente con el segundo?}



The literature on left-behind places (LBPs) provides a powerful analytical and normative framework to understand why some territories remain trapped in low-opportunity trajectories and to design interventions tailored to their specific constraints \citep{rodriguez2024geography}. The LBP concept  \citep{pike_rachel_2023left,mackinnon_festival,faggian2024three} describes areas that have been sidelined by dominant trajectories of economic transformation and territorial governance, characterized by the accumulation of interrelated disadvantages \citep{davenport2020levelling, hendrickson_2018}. This framework is multidimensional, integrating demographic (e.g., weak population dynamics and age imbalances), economic (e.g., low dynamism, weak entrepreneurship, and chronic underinvestment), and accessibility/connectivity dimensions (high temporal distance to markets, limited services and networks). In essence, the concept transcends purely economic metrics to encompass socio-political exclusion, institutional neglect, and a pervasive feeling of peripheralization and institutional disengagement \citep{dijkstra2020geography}.\footnote{These complex conditions can generate disaffection, shape political preferences, and intensify territorial tensions \citep{tomaney2021levelling}.}


{\color{red} (El párrafo que viene ahora es larguísimo)}

Recent studies have conceptualised left-behindness as a multidimensional and dynamic process driven by changes in multiple social and economic factors.\footnote{For instance, \cite{connor2024gets} define left-behindness based on the evolution of poverty rates, median household income, unemployment and educational attainment, measuring it through changes in a rank-based composite index. \cite{Perancho2025} delve deeper into this perspective, constructing a multidimensional synthetic index to identify left-behind municipalities (LAU2) across Europe. \color{red} (quizás esta footnote sea prescindible) } Our study addresses this limitation by incorporating multiple dimensions of analysis.\footnote{For instance, in Spain, this phenomenon is primarily associated with pronounced demographic decline, contrasting with the UK's historical focus on deindustrialized towns.} Despite the strengths of this integrated conceptual framework, the empirical literature on LBPs still presents several major methodological limitations, some of which this study seeks to address. First, a significant methodological weakness is the reliance on highly aggregated territorial units (e.g., NUTS2 or NUTS3 in European terminology). While empirical work has often targeted specific regional archetypes, such as post-industrial or peripheral rural zones \citep[see][]{guilluy2019twilight, mackinnon2021reframing}, this macro-scale approach obscures the local heterogeneity of disadvantage. Therefore, the characteristics that define LBPs must be localized at the territorial level, which leads to an aggregation bias in regional analyses. Following this critical insight, and validated by recent local analysis \citep{Perancho2025,alama2021drivers}, we adopt a more disaggregated territorial level, namely, municipalities, corresponding to LAU2 in European Union terminology.\footnote{LAU2 refers to the Local Administrative Units level 2, which also corresponds to the municipal in most other European Union countries.} Second, despite growing recognition in contemporary approaches that geography acts as an active component shaping economic outcomes {\color{red} (De la Roca et al., 2024) (missing citation)}, this perspective remains far from standard in empirical analyses. In particular, the role of spatial context, captured through second-nature factors such as accessibility and market potential, remains insufficiently integrated into the identification of left-behind places. Ignoring these dimensions may lead to incomplete or even misleading diagnoses of territorial socio-economic conditions \citep{iammarino2019regional}. Our study addresses this gap by explicitly incorporating geographic factors into the identification of left-behind municipalities.  Third; according to \cite{phelps2025abandoned}, \textit{``increased attention will be needed to the systemic, multiscalar, and relational processes that produce territorial abandonment, decline, and disposability. This calls for moving beyond symptom-based typologies toward place-sensitive frameworks that interrogate the political, economic, and social mechanisms underpinning the processes that generate left-behindness''} \citep[p.8]{phelps2025abandoned}. Building on this perspective, we explore how the variables considered are associated with the probability of belonging to different diagnosed categories of left-behind places (LBP).




In the Spanish context, on which we focus, a critical research gap remains: while existing studies explore related processes such as depopulation or territorial vulnerability \citep{gonzalez2023understanding}, comprehensive and data-driven typologies focusing on the diverse municipal reality of Spain remain scarce.\footnote{Recent work by \cite{Perancho2025} represents an important step forward by adopting a multidimensional and municipal-level framework to identify left-behind municipalities (LAU2) across Europe. However, despite its multidimensional nature, this approach is based on the construction of a synthetic index that condenses multiple territorial characteristics into a single measure. While this provides a useful overall ranking of municipalities according to their level of disadvantage, it does not allow for the identification of distinct patterns or profiles of left-behindness. Consequently, important differences across municipalities remain hidden, limiting the scope for designing differentiated and place-based policy responses.} This shortcoming motivates the central research questions of this study: what are the main types or forms of municipal left-behindness in Spain?, does second-nature geography matters? and how the different variables influence the probability of belonging to these forms of left-behindnes in the Spanish municipalities.



Our paper investigates left-behind places (LBPs) in Spain with the aim of explaining both municipalities' proximity to vulnerability thresholds and their likelihood of falling into specific LBP categories. This central aim requires, as a first step, to offer a data-driven typology of Spanish municipalities according to both the degree and the nature of structural disadvantage. Recent empirical studies have increasingly applied cluster analysis to identify typologies of territorial disadvantage and left-behindness.\footnote{For example, \cite{le2025identifying} and \cite{velthuis2024regional} use $k$-means clustering to identify distinct groups of NUTS3 regions in the EU15 based on multiple economic, demographic and social indicators. Similarly, \cite{hedlund2016mapping} apply cluster methods to map the socioeconomic landscape of rural Sweden and to derive a typology of rural areas, while \cite{jessen2024role} use clustering techniques to identify different variants of left-behind places in Denmark. Earlier contributions in regional economics, such as \cite{kronthaler2005economic}, also applied cluster analysis based on regional growth factors to classify East and West German regions according to their economic capability.} Most of these cluster-based studies on territorial disadvantage implicitly treat geography as a contextual factor that helps explain differences between clusters, rather than as a constitutive dimension of disadvantage itself. However, in the case of left-behind places, disadvantage often reflects not only internal socioeconomic and demographic conditions, but also municipalities' functional position within wider territorial systems. Ignoring this dimension may therefore overlook forms of left-behindness that are inherently territorial in nature. Hence, a contribution of our study builds on this insight by examining the effects of the incorporation of second-nature geographic factors into cluster-based typologies. We argue that the consideration of these dimensions is fundamental for understanding the structural conditions that differentiate dynamic municipalities from those facing persistent decline \citep{velthuis2024regional}. In addition, because of working at the municipal scale, our approach captures territorial heterogeneity in much greater detail, identifying distinct spatial patterns in economic resilience, labour markets, and infrastructure.

{\color{red} (Contactar con los autores que hacen cluster analysis)}


Our results reveal that, while the overall cluster structure remains broadly stable, the explicit incorporation of second-nature geographic variables brings to light a significant increase in vulnerability across a substantial number of municipalities. Rather than simply preserving initial classifications, the consideration of geography leads to a non-negligible share of municipalities being reassigned towards more vulnerable profiles. These reassignments are not random, but are particularly concentrated among municipalities with intermediate characteristics, suggesting that geographic constraints play a critical role in pushing these territories towards less favourable positions within the cluster structure.



Moreover, post-clustering sensitivity analysis indicates that second-nature geographic variables do not merely operate at the margins of the classification, but actively contribute to redefining the relative proximity between clusters for a meaningful subset of observations. Consistent with this view, geography does not just refine the classification, it reveals latent structural disadvantages that were previously less visible. Rather than acting as a neutral adjustment, the explicit inclusion of these variables underscores the extent to which territorial vulnerability is shaped by spatial constraints, reinforcing the idea that many municipalities are more vulnerable than initially suggested when geography is not fully accounted for. Consequently, a central contribution of the paper is to highlight the role of second-nature geography in shaping these processes. Variables capturing accessibility to urban and semiurban centres and market potential emerge as critical factors in explaining municipalities’ position relative to vulnerability frontiers and their probability of belonging to specific LBP types. From a policy perspective, the findings emphasize the relevance of the municipal scale and support a shift towards more integrated, systemic place-based policies that explicitly account for multi-scalar dynamics and the embeddedness of local trajectories within broader regional systems.





The study proceeds as follows. After this introduction, Section \ref{sec:theoretical_background} reviews the relevant literature and provides the theoretical background. Section \ref{sec:territorial_structure} describes the territorial framework and data used, while Section \ref{sec:methodology} presents the methodology, and Section \ref{sec:results} the results. Results are further validated and analysed more deeply in Section \ref{sec:post_cluster} and, finally, Section \ref{sec:conclusions} outlines the main conclusions.



%This analysis contributes to the broader debate on the design of place-sensitive territorial policies. Moving beyond one-size-fits-all approaches, our typology supports the development of more targeted interventions adapted to the specific needs and vulnerabilities of different types of left behind places in Spain. Understanding this diversity is essential for fostering inclusive and territorially balanced development strategies.





\section{Theoretical background}
\label{sec:theoretical_background}





\subsection{Geography, clusters and left-behindness}
\label{sec:geography_clusters}


Since the 2008 crisis, the concept of ‘left behind places’ has gained prominence as a key lens through which to understand growing geographical inequalities \citep{pike2024left}. Although the term encompasses a wide range of political, economic, social and cultural dimensions of territorial disadvantage, these manifestations are ultimately rooted in spatial processes. From a place-based perspective, it corresponds to areas that bring together a number of specific geographic characteristics (at different scales) that have been left behind by unequal patterns of economic transformation \citep{martin2021levelling}, highlighting the central role of location in the emergence and persistence of these inequalities. While \cite{pike2024left} outline the various dimensions that characterize left-behind places (including economic, social, environmental, political, institutional, and cultural aspects), the geographical dimension is often notably absent or treated implicitly. However, elements such as proximity to semi-urban and urban centres, and the territorial connectedness of each area, play a fundamental role in shaping local development dynamics and, ideally, should be explicitly modelled as structural drivers of left-behindness. In this regard, the importance of geography as a key driver of local development is supported by an extensive literature analyzing the role of territorial connectivity and linkages between areas with different levels of urbanization. Concepts such as spatial spillovers, backwash effects, and spread effects are widely used to describe these dynamics.\footnote{These dynamics are closely linked to the generation and diffusion of agglomeration economies associated with the spatial concentration of population and employment. Spatial proximity gives rise to positive externalities that may spill over to neighbouring areas, implying a simultaneity between local and interregional effects. In this context, connections with more dynamic regions can generate spread effects that benefit adjacent areas facing greater structural challenges. At the micro level, these externalities may stem from Marshallian mechanisms (Duranton and Puga, 2004), as well as from pecuniary externalities associated with increasing returns and transport costs (Ottaviano and Thisse, 2001). \color{red} (poner las citas como toca)} {\color{red} (see Alamá et al., 2025c) (citar bien) }


{\color{red} (No sé si en este párrafo estamos repitiendo cosas de la intro; es muy típico de la IA, y ya me lo han dicho algunos editores amigos).}



%This body of work, often relying on methodologies based on the system of simultaneous equations \citep{carlino1987determinants}, measures the effect of accessibility on municipal development through the estimated elasticity of local population with respect to employment in surrounding areas.

%More recently, some contributions have factored spatial and sectoral heterogeneity into the analysis, considering that the indirect effect of interaction depends on the types of territory involved \citep{feser2006harnessing, bosworth2018economic}. In the Spanish context, authors like \cite{alama2025spatio} analyse the full set of territorial interactions within regional systems by explicitly modelling the population–employment relationship, underscoring the importance of spatio-sectoral heterogeneity in shaping interrelations among local areas. .{\color{purple}


A recurrent methodological issue in the cluster-based literature {\color{red}(citas de la intro)} on territorial disadvantage concerns whether the inclusion of geography introduces distortion or, instead, enhances explanatory power. Many studies deliberately exclude geographic variables from the clustering procedure, focusing instead on demographic and socioeconomic indicators to identify territorial typologies \citep[see, for instance,][]{hedlund2016mapping, jessen2024role, velthuis2024regional, kronthaler2005economic}. This choice is commonly justified by the concern that geographical distance, location or contiguity may distort clustering outcomes. Specifically, geographic variables tend to exhibit strong spatial heterogeneity \citep{anselin1988spatial}. When incorporated into clustering algorithms, this may result in groupings driven primarily by spatial proximity, with neighboring municipalities being assigned to the same cluster regardless of their underlying socioeconomic structures. In such cases, clusters may resemble broad regional patterns (e.g. north--south or coastal--inland divides) rather than distinct configurations of territorial disadvantage. A related rationale is the intention to identify ``pure'' socioeconomic typologies, whereby municipalities are first classified exclusively on the basis of demographic and socioeconomic characteristics and geographical factors are subsequently analyzed \textsl{ex post}. This strategy is internally consistent when the research question seeks to abstract from spatial context. However, recent contributions have increasingly emphasized that such an abstraction may be problematic in the context of left-behind places. In particular, \citet{jessen2024role} argues for the need to incorporate greater ``geographical wisdom'' when identifying left-behind regions, highlighting that territorial disadvantage is deeply embedded in region-specific spatio-temporal contexts. While this perspective underscores the importance of geography for understanding left-behindness, geographic factors are still largely treated as contextual or interpretative elements, rather than being directly incorporated into the clustering procedure itself.



This study builds on this insight: in the context of left-behind places, second-nature geography is not merely a descriptive backdrop, but a structural dimension shaping local development trajectories.  Hence, our study adopts a different methodological approach based on cluster analysis.\footnote{This technique has been proved to be useful when classifying territories according to their  structural characteristics \citep[see, for instance][]{ebbekink2013s}, as a way to identify patterns conducive to greater territorial development.} This approach departs from much of the existing literature on left-behind places, which typically relies on cluster-based typologies that abstract from geographical factors \citep{le2025identifying, velthuis2024regional, hedlund2016mapping, jessen2024role, kronthaler2005economic}. Second-nature geography is not introduced to explain clusters ex post, but to allow the clustering procedure itself to reflect functional territorial disadvantage. The resulting classifications can facilitate the design of policy interventions tailored to the specific characteristics of each     .\footnote{Although from a conceptual perspective the term ``cluster policy'' has gained importance in the field \citep{brakman2013reflections}, methodologically it has been considered less frequently.} Thus, our study assesses how municipal vulnerability is shaped by demographic, socioeconomic, and geographic factors, with particular attention to second-nature geographic indicators as drivers of territorial disadvantage.




\subsection{Analytical strategy}
\label{sec:analytical_strategy}


To assess the role played by geography in the identification of left-behind municipalities, we adopt a two-scenario clustering approach. This comparative strategy allows us to examine how the inclusion of second-nature geographic indicators affects the classification of municipalities and to identify cases in which municipalities are reassigned once spatial constraints are taken into account. In particular, it enables the identification of municipalities whose classification as left-behind or non-left-behind changes when geographical factors are incorporated. Specifically, we define a comprehensive set of demographic, socioeconomic, and geographic variables to capture the diversity of territorial contexts.





To evaluate the contribution of the geographical dimension to the classification, we construct two comparative scenarios: one based exclusively on demographic and socioeconomic variables, and another one that additionally incorporates geographical indicators. In this regard, in Scenario I, our cluster analysis is based exclusively on demographic and socioeconomic variables---in line with much of the existing literature \citep{velthuis2024regional}. Under Scenario II, the cluster analysis incorporates second-nature geographic variables directly into the clustering process. This comparative approach enables us to answer two central questions: (i) to what extent does geography influence the classification of municipalities?; and (ii) to what degree can such geographic variables tip the balance, pushing certain areas towards less favourable development trajectories and thus leading them to be classified as left-behind places?




Distance to urban or semi-urban cities and market potential are introduced as geographic indicators. In line with the distinction between first and second nature geography, we classify market potential as a geographic control, grounded in the spatial configuration of locations and distance-related frictions. At the same time, it is important to acknowledge that market potential embodies second nature geography, as it reflects endogenous spatial-economic processes arising from the concentration of demand and economic activity. While market potential is geographic in origin, its informational content is hybrid, capturing how spatial positioning translates into economically meaningful access to markets through circular and cumulative mechanisms central to economic geography nature \citep{bathelt2024nature}. This interpretation aligns with contemporary approaches that conceptualize geography not as a passive backdrop but as an active component shaping economic outcomes \citep{de2024journal}. Distance to urban or semi-urban cities is likewise classified as a geographic indicator, as it directly captures the spatial positioning of a location relative to urban centers. This measure reflects exogenous spatial accessibility and transport-related frictions that condition interaction costs and connectivity across space. Unlike market potential, distance does not embed endogenous economic outcomes, but rather represents a foundational geographic constraint through which second nature processes may subsequently operate \citep{jessen2024role}. As for the geographic indicators employed here, such as distance to urban and semi-urban centres and measures of market potential, they aim to capture functional accessibility rather than mere locational attributes, reflecting the extent to which municipalities are connected to economic opportunities and embedded within broader territorial systems; hence, they complement purely demographic and socioeconomic characteristics.\footnote{Unlike simple measures of Euclidean distance, the distance indicator used in this study (see Equation \eqref{eq:eq9} below) combines distances to urban and semi-urban centres through a non-linear formulation that captures the mitigating effect of access to at least one type of centre on spatial isolation. As a result, municipalities that are geographically close may exhibit substantially different distance values depending on their functional connectivity. This construction helps to limit the risk of strong spatial autocorrelation typically associated with purely geographical distance measures, as the indicator does not capture location \textsl{per se}, but relative access to economic and service centres. While some degree of spatial dependence is inevitable in territorial data, this approach reduces the likelihood that clustering outcomes are driven mechanically by spatial proximity rather than by structurally meaningful differences in accessibility. {\color{red} (quizás esta nota es prescindible, o se pueda reducir)}}

% According to our results, and as we will see below, when distance to urban/semi-urban centres and market potential are incorporated into the analysis, nearly one-fifth of municipalities shift classification, showing that geographic disadvantage is a fundamental, independent driver of territorial marginalization that demands explicit policy attention.



{\color{red} Avoid repetition; decirle que trate de escribir como Andrés.}


\section{Territorial framework and data}
\label{sec:territorial_structure}

\subsection{Territorial framework}
\label{sec:territorial_framework}


Spain's administrative classification is built upon a multi-level structure composed of 17 autonomous communities (\textsl{Comunidades Autónomas}, CCAA) and 50 provinces, which define the regional and intermediate governance tiers. The most fundamental unit of local administration, and the primary focus of this research, is the municipality (LAU1). In Spain there are  approximately 8,131 municipalities. The distribution of these units across the regional administrative boundaries (autonomous communities) is represented in Figure \ref{fig:map_ccaa_municipal_labels}.\footnote{Spain's Autonomous Communities  correspond to NUTS2 regions in the European Union's Nomenclature of Territorial Units for Statistics. Established through the 1978 Constitution and subsequent autonomy statutes, Spain's 17 autonomous communities (plus 2 autonomous cities) represent the primary level of sub-national governance, exercising legislative and executive authority in devolved policy areas. This decentralised structure creates heterogeneous institutional and policy frameworks across regions, making within-region comparisons particularly relevant for understanding municipal-level dynamics.}

{\color{red} Comentar lo que es en European terminology; citar algún estudio de Dani Tirado; ojo, municipalities son LAU2.}



The {\color{red}LAU1 (ojo, municipalities son LAU2)} level exhibits marked heterogeneity across key dimensions such as population size, demographic structure, and socioeconomic profile. To capture this variation, {\color{red}Eurostat (2021) (citar)} classifies Spanish municipalities according to their degree of urbanization (Figure \ref{fig:urbanizacion_municipios}).\footnote{Following Eurostat’s standard methodology, we apply the DEGURBA framework {\color{red}Eurostat (2021) (citar)}, which distinguishes three categories: urban centres, intermediate density areas, and predominantly rural areas.} The analysis highlights the highly unbalanced nature of Spain's local landscape: the vast majority of municipalities (83.50\%) are predominantly rural, yet they account for only 13.23\% of the population in 2022. In contrast, intermediate density areas represent 32.12\% of municipalities and concentrate 38.16\% of the population, while urban centers, despite representing barely 3.10\% of municipalities, have a population share of 54.65\%.





\subsection{Data sources}
\label{sec:data}


\subsubsection{Demographic variables}
\label{sec:demographic_variables}

%regpopgrowccaa------------------------------------------------------------

In this section, to capture the demographic dimension, we describe the variables used to characterize differences in population trends across Spanish municipalities. These variables offer key insights into the underlying demographic dynamics that drive population growth or decline. By analyzing factors such as age distribution, migration patterns, and population change, we aim to identify structural differences among municipalities and their implications for long-term demographic evolution \citep{jessen2024role}. The cluster analysis and the classification of municipalities are based on data for year 2022 and, to assess long-term population growth in each municipality relative to its regional context, we calculate growth rates using the year 2000 as a fixed baseline.




In Equation \eqref{eq:eq1} we represent population growth ($\Delta POP_i$) in municipalities relative to regional population growth over the period 2000--2022, using year 2000 as baseline. If a municipality's population growth matches that of its home region (autonomous community), the indicator equals zero. A positive value indicates that the municipality's population has grown at a faster rate than its home region's population, while a negative value reflects the opposite.

{\color{red} (Mapa de Paco Goerlich de Linkedin).}



\begin{equation}
\label{eq:eq1}
\Delta POP_i = \left(\frac{POP_i^{t}}{POP_i^{t-n}}\right) - \left(\frac{POP_{ccaa(i)}^{t}}{POP_{ccaa(i)}^{t-n}} \right)
\end{equation}

\noindent where $\Delta POP_i$ is the population growth of municipality $i$ relative to its home region's population growth between years 2000 and 2022, $POP_i$ is municipality $i$'s population, and $POP_{ccaa(i)}$ is its home region's population.


%-----------------------------------------

Demographic variables at the municipal level in Spain include key indicators such as the elderly dependency ratio ($DEPEND$), defined as the ratio of the population aged 65 and over to the working-age population, and the share of the foreign population under 16 years of age ($ShFPOP16$). These measures are formally defined in Equations \eqref{eq:eq2} and \eqref{eq:eq3}:



\begin{equation}
\label{eq:eq2}
DEPEND_i =
\frac{POP65_i}{POP1664_i}
\end{equation}

\begin{equation}
\label{eq:eq3}
ShFPOP16_i=
\frac{FPOP16_i}{FPOP_i}
\end{equation}


In Equation \eqref{eq:eq2}, the ratio $DEPEND_i$ corresponds to the old-age dependency index in municipality $i$, $POP65_i$ is the total population aged 65 years or older, and $POP1664_i$ is the working-age population (aged 16 to 64). As for Equation \eqref{eq:eq3}, $ShFPOP16_i$ measures the share of the foreign population under 16 years of age in municipality $i$, $FPOP16_i$ represents the total foreign population under the age of 16, and $FPOP_i$ is the municipality's total foreign population.




High values of the elderly dependency index in Equation \eqref{eq:eq2} characterize municipalities experiencing significant demographic aging and low generational renewal. Likewise, the share of the foreign population under 16 years old in Equation \eqref{eq:eq3} reflects the presence of young immigrant groups that can partially mitigate aging and contribute to demographic regeneration in some municipalities. A low share of foreign youth might indicate a lack of demographic dynamism, whereas higher shares may be associated with municipalities that, despite structural fragilities, remain attractive to certain population groups \citep{kuhn2015peripheralization,leibert2016peripheralisation}.




As for net migration ($NETMIGR_i$), it plays a significant role in the demographic revitalization of municipalities, particularly in inland areas facing persistent depopulation processes:\footnote{In these territories, characterized by negative natural growth due to population ageing and low fertility rates, immigration flows are essential to offset population losses stemming from natural decrease. Empirical evidence shows that, in the absence of positive net migration, these municipalities would experience a much sharper demographic decline, threatening their long-term social and economic sustainability \citep[see][]{AlamaSabaterBernatBudi2024}.}



\begin{equation}
\label{eq:eq4}
NETMIGR_i = \frac{IMMGR_i - EMIGR_i}{POP_i}
\end{equation}

\noindent where $IMMGR_i$ denotes the inflow of migrants to municipality $i$ during the reference period, $EMIGR_i$ the outflow (emigrants), and $POP_i$ the municipality's total population.


Furthermore, we consider population density as a key demographic and geographic factor reflecting the spatial distribution of inhabitants within municipalities, being calculated as the ratio of population ($POP_i$) to municipality $i$'s total surface ($SURFACE_i$) (in square kilometers):



\begin{equation}
    \label{eq:eq11}
    DENSPOP_i = \frac{POP_i}{SURFACE_i}
\end{equation}



\subsubsection{Economic variables}
\label{sec:economic_variables}


Economic variables are essential to capture structural differences across municipalities and to classify them according to their level of economic disadvantage. In this sense, we use income per capita ($INCOMEpc_i$) as an indicator of the regional level of economic development \citep[see][]{velthuis2024regional,fiorentino2024left,Perancho2025}.\footnote{We intentionally avoid rescaling $INCOMEpc_i$ relative to the home region's average. Doing so would distort the interpretation of economic disadvantage in the clustering process. Likewise, we do not rescale income by the national average, as this transformation would be redundant. \color{red} (esta nota se podría quitar)}







With regard to sectoral patterns, left-behind places are typically characterized by a weak economic structure, often marked by limited diversification and a strong dependence on traditional sectors. In many cases, their economy relies heavily on agriculture which, despite its relevance, tends to generate low value-added employment and remains highly vulnerable to external shocks such as climate change and market fluctuations. In some territories, tourism also plays a significant role; however, its predominantly seasonal nature often results in unstable employment patterns and constrains long-term economic sustainability \citep{fiorentino2024left,jessen2024role}.

Against this background, analyzing the economic structure of left-behind places requires close attention to local employment patterns, specifically the distribution of employment across key economic sectors. These territories are frequently characterized by limited economic diversification, with a substantial share of employment concentrated in traditional activities. Consequently, incorporating sectoral employment shares into the clustering framework enables capturing fundamental differences in local economic models and development trajectories, and to distinguish between different types of economic lag among municipalities. Our employment data are sourced from the contract register provided by municipalities through the Public Employment Service (SEPE).


To capture the sectoral specialization of municipalities, we compute the share of contracts in each economic sector relative to the total number registered within each municipality. Specifically, we focus on two sectors that are particularly relevant in the context of left-behind places: industry ($ShINDUST$) and agriculture ($ShAGRICULT$):




\begin{equation}
\label{eq:eq6}
ShINDUST_i = \frac{INDUST_i}{TOTALCONTR_i}
\end{equation}

\begin{equation}
\label{eq:eq7}
ShAGRICULT_i = \frac{AGRICULT_i}{TOTALCONTR_i}
\end{equation}



The presence of bank branches also plays a crucial role in assessing different types of left-behind places and identifying spatial clusters of economic disadvantage {\color{red} (paper Flogel)}. Access to financial services is a key determinant of local economic vitality, as it facilitates business investment, supports household financial stability, and enables public and private initiatives that drive regional development. In many left-behind areas, a declining number of bank branches signals financial exclusion, limiting access to credit, savings, and other essential financial services {\color{red} (refs-banking deserts, all recent citations Mari Carmen, Rodrigo, etc.)}. This can particularly affect small businesses, entrepreneurs, and rural communities that rely on local banking infrastructure for their economic activities {\color{red} (paper Abedifar Ir\'an)}. The absence or reduction of branches may also reflect broader structural challenges, such as depopulation, aging demographics, and weak economic performance, reinforcing a cycle of economic decline {\color{red} (citations)}. Furthermore, analyzing the distribution of bank branches helps to differentiate between various types of left-behind places. Some areas may suffer from economic stagnation but retain financial services due to tourism or industrial activity, while others, particularly rural or remote regions, may experience both economic decline and financial exclusion. Hence, identifying clusters of municipalities with low banking access can highlight regions where intervention is needed to promote financial inclusion and stimulate local economic recovery. In this regard, the integration of bank branch data into the study of left-behind places which, up to now, {\color{red} only} \cite{Flogel} has done to some extent, provides valuable insights into regional disparities, the dynamics of economic marginalization, and the potential for policy interventions aimed at fostering economic resilience and development.



%2020–2025? or 2015-2022?
In our context, we include the share of bank branches relative to the population in each municipality during the period 2020--2025. This enables capturing trends in financial accessibility and assess their impact on local economic conditions, helping to identify areas experiencing financial exclusion and those maintaining or improving access to banking services:
%
\begin{equation}
\label{eq:eq8}
ShBRANCHES_i = \frac{BRANCHES_i}{POP_i}
\end{equation}




\subsubsection{Second-nature geographic indicators}
\label{sec:geographic_distance_variables}

{\color{red} (quizás se podría poner un párrafo introductorio a second nature, avoiding repetition)}

We include an explicit analysis of \underline{accessibility} by measuring the distance in kilometres from rural municipalities to semi-urban and urban centres. This approach allows us to assess the degree of spatial isolation and the challenges residents face in accessing essential services, economic opportunities and financial institutions. To better capture accessibility effects, we define a composite distance measure that reflects proximity to either type of centre. Specifically, we construct a new distance function as follows:

\begin{equation}
\label{eq:eq9}
DISTANCE_i = \frac{URBANDIST_i \times SEMIDIST_i}{URBANDIST_i + SEMIDIST_i}
\end{equation}

{\color{red} (habría que explicar qué es cada cosa en la fórmula)}

This formulation emphasizes the mitigating effect that proximity to at least one centre (urban or semi-urban) has on spatial isolation. The index takes lower values when a municipality is close to either type of center, thus reflecting reduced exposure to isolation-related constraints, while higher values are observed only when municipalities are distant from both urban and semi-urban centres, indicating greater remoteness and potentially higher vulnerability.

Importantly, this non-linear construction captures functional accessibility rather than absolute spatial proximity. Two geographically close municipalities may therefore display markedly different distance values depending on their access to urban or semi-urban centres. This feature helps to limit the risk of strong spatial autocorrelation typically associated with simple Euclidean distance measures, as the indicator does not capture location per se, but relative access to economic and service hubs. While some degree of spatial dependence is inevitable in territorial data, this approach reduces the likelihood that accessibility effects mechanically mirror spatial proximity. {\color{red} (quizás este párrafo podría ir a footnote)}


Understanding these spatial constraints is essential for identifying areas with limited functional connectivity, where targeted place-based policies aimed at improving mobility, infrastructure and access to services may enhance territorial integration and development prospects \citep{paez2012measuring,Mustafa2025}.


Accessibility is also determined by considering the characteristics of neighboring municipalities, as the surrounding area can be regarded as a potential employment zone or territorial unit for each municipality. In this context, to emphasize the relevance of neighboring characteristics in the cluster analysis, we have included the market potential indicator as a geographical and distance-based variable. In the context of New Economic Geography, ``market potential'' captures the attractiveness of a specific location for economic activity \citep{krugman1991increasing,harris1954market}. Thus, market potential is computed using a distance-weighted measure of surrounding economic activity. Specifically, for each municipality, we calculate the contribution of neighbouring municipalities as a function of their income levels, weighted by the inverse squared distance between centroids. Distances are computed as Euclidean distances between municipal coordinates, and self-interactions (zero distance) are excluded from the calculation. Formally, the market potential indicator assigns greater weight to nearby municipalities, reflecting the rapid decay of economic interactions with distance, and is defined using an inverse-distance-squared weighting scheme. For a given municipality $i$, its market potential at time $t$ can be calculated as follows:

\begin{equation}
    \label{eq:eq10}
    PM_{it} = \sum_{j=1}^{n} \frac{Income_{j}}{d_{ij}^2}
\end{equation}

\noindent where $Income_{j}$ represents the income per capita of municipality $j$, $d_{ij}$ represents the distance between municipality $i$ and municipality $j$, and $n$ is the total number of municipalities


All variables introduced in this section are described with more precision in Table \ref{tbl:1}, which not only provides definitions and descriptions of each of the variables used, but also the statistical sources for each of the variables.

{\color{red} (Comprobar en la tabla que las expresiones corresponden a las ecuaciones)}




\section{Methodology and clustering process}
\label{sec:methodology}



The identification of left-behind places poses a significant methodological challenge due to the multidimensional nature of territorial disadvantage and the high degree of heterogeneity across local contexts. In countries such as Spain, where more than 8,000 municipalities coexist with markedly different demographic, economic and accessibility profiles, approaches based on a single indicator or on \textsl{a priori} classifications risk oversimplifying complex spatial realities. Capturing these diverse development trajectories therefore requires methods capable of jointly considering multiple dimensions of disadvantage without imposing restrictive assumptions on their relative importance. In this context, cluster analysis offers a particularly suitable analytical framework.\footnote{Cluster analysis is a fundamental technique in exploratory data analysis aimed at identifying natural groupings in complex datasets \citep{johnson2002applied}. It seeks to maximize within-group similarity while ensuring between-group dissimilarity. This data-driven approach is especially advantageous in settings with a large number of observational units, such as the Spanish municipal system, as it allows municipalities to be grouped according to similarities across a set of relevant characteristics, while letting the data reveal endogenous patterns of territorial differentiation. \color{red} (no sé si vale la pena ponerlo)}  When applied to municipal-level information, cluster analysis enables the identification of distinct types of left-behind places, reflecting different combinations of economic structure, socioeconomic conditions, demographic dynamics and accessibility constraints.

\color{black}






Several methods can be employed to perform cluster analysis, each with its own assumptions and computational approaches. Partitioning algorithms divide the set of objects into a predetermined number of groups by minimizing within-cluster variation, while hierarchical algorithms construct a nested structure of clusters either by successively agglomerating observations (bottom-up) or by recursively splitting them (top-down). Additionally, Gaussian mixture models offer an alternative probabilistic approach, assuming that the data are generated from a mixture of Gaussian distributions and estimating their parameters to infer cluster memberships.




%This section outlines the methodological approach adopted in this study for clustering Spanish municipalities, discussing the main clustering techniques available and justifying the choice of method used.

%Cluster analysis is a family of statistical techniques that aim to classify a set of individuals---in this case, Spanish municipalities---into groups (or clusters) such that the members within each group are as similar as possible to one another, and as dissimilar as possible from those in other groups.


In our case, and given the large number of municipalities and the high level of data granularity, partitioning methods such as \emph{k}-means are particularly suitable, as they allow for an efficient and scalable classification of territorial units while yielding clusters that are easy to interpret from a regional development perspective \citep{hartigan1979algorithm}, \citep{jain2010data}.
%In our case, partitioning methods could be the most appropriate choice for classifying Spanish municipalities, due to their advantages for handling large and granular datasets, offering a scalable and efficient solution along with a straightforward interpretation of the resulting clusters {\color{red} (refs)}.
In contrast, hierarchical methods tend to be computationally intensive and less appropriate for large samples, while Gaussian mixture models, although more flexible, rely on strong distributional assumptions and can be sensitive to initialization and convergence issues. Moreover, partitioning techniques are widely used in similar empirical studies focused on territorial or municipal classifications \citep{legendre2012clustering}, \citep{horn1995solution}.

In partitioning algorithms, the similarity between municipalities is quantified through a distance metric. Since all variables in the dataset are numerical, three alternative distance measures were considered: Euclidean, Mahalanobis, and Canberra. Each of these metrics relies on different underlying principles and presents specific advantages depending on the structure of the data. In this study, the Euclidean distance is adopted as the main metric for the clustering procedure. A detailed description of the alternative distance measures considered and their properties is provided in Appendix \ref{sec:appendix_A}.


%After testing the three metrics, we selected Euclidean distance for the clustering process, as it best aligned with our goal of grouping municipalities by overall similarity in socioeconomic indicators. While Canberra distance was less effective due to the homogeneous nature of our variables, Mahalanobis distance accounted for correlations but did not suit our conceptual focus on absolute proximity in a multidimensional space. A detailed explanation can be found in Annex A.




%- dades faltants:
%la renda mitjana per càpita municipal l’he haguda d’obtenir de les bases de dades experimentals de l’INE, a causa del nivell elevat de desagregació municipal. Això implica que faltarien moltes dades. Afortunadament, aquest problema només afectava la variable de la renda municipal, fet que em permetia dissenyar un algoritme d’interpolació per aproximar els valors faltants. Aquest es basa en una mitjana d'una aproximacio espacial i una temporal (per obtindre un primer valor d'acord a unes condicions semblants a la del municipi, i després adaptar-lo a la tendencia d'aquest)
%Per a resoldre les poques dades faltants que quedaven hi havien dos solucions, suprimir els municipis de la base de dades o imputar aquests valors amb l'instrucció 'imputePCA' de la llibreria 'missMDA' de R. Com que hi faltaven poques dades i no afectaria als resultats finals vaig emprar imputePCA.


Partitioning methods divide a dataset composed of \(n\) municipalities into \(K\) clusters, with \(K<n\). The number of clusters must be specified in advance. In this study, the optimal value of \(K\) is determined by combining two of the most commonly used criteria in the clustering literature: the elbow method and the maximization of the average silhouette index\footnote{A detailed description of the elbow method and the silhouette index is provided in Appendix \ref{sec:appendix_B}.} \citep{januzaj2023determining}. Although the elbow method is commonly sufficient in the literature to determine the optimal number of clusters, in our case it did not provide a clear inflection point, as shown in Appendix \ref{sec:appendix_B}. For this reason, the final selection of \(K\) relies on the maximization of the average silhouette index.





The standard partitioning method using Euclidean distance is the \emph{k}-means algorithm.\footnote{see Appendix \ref{sec:appendix_c}} The choice of the $k$-means algorithm is also motivated by the characteristics of the dataset used in this study. The analysis relies on a large cross-sectional dataset covering thousands of municipalities, which exhibit substantial heterogeneity in demographic, economic, and geographic conditions. Against this backdrop, partitioning methods such as $k$-means are particularly suitable, as they allow the identification of relatively homogeneous groups of municipalities while remaining computationally efficient for large samples.








%Once the clustering procedures have been applied, it is essential to evaluate the quality of the resulting classifications. Among the available validation tools, we use the silhouette coefficient, a widely accepted internal measure that assesses how well each observation has been classified within its assigned cluster \citep{januzaj2023determining}. A detailed explanation of the silhouette indicator is provided in Appendix \ref{sec:appendix_B}.



%In our study, we calculated the silhouette coefficient for each clustering configuration generated by the partitioning methods. The final classification was selected as the one that maximized the average silhouette value \(\bar{s}\) in Equation \eqref{eq:eq28} of Appendix \ref{sec:appendix_B}.


\begin{comment}


\subsection{Clustering process}
\label{sec:clustering_process}


The classification of Spanish municipalities into distinct structural groups was performed using the $k$-means clustering approach, a widely accepted partitioning algorithm suitable for large datasets, as explained above.

The objective is to construct a robust and meaningful typology of municipalities based on relevant structural characteristics, thereby enabling the identification of the heterogeneous profiles of Left-Behind Places (LBPs).



To ensure the robustness of the classification and to isolate the specific role of geographical factors, the clustering was executed in two comparative {\color{purple} scenarios}. This {\color{purple}comparative} design serves as a relevant analytical experiment to detect any ``hidden'' marginality trait that might be obscured when spatial or {\color{purple} geographical} variables {\color{red} (spatial or all the variables?. SPATIAL)} are not explicitly considered.




{\color{purple}The two scenarios  enables us to assess} the role played by distance-related  factors in the classification process and detects whether excluding geographic variables leads to a typology that underestimates the structural disadvantage of municipalities with limited accessibility. By comparing both typologies, one based solely on demographic and socioeconomic indicators and another that also incorporates geographical and distance-related factors, the analysis aims to assess the extent to which geographic isolation contributes to a municipality's vulnerability, and to what degree it may remain overlooked if spatial factors are not explicitly considered. These two stages are defined as follows:




\begin{enumerate}
    \item \textbf{Scenario I:} Clustering was performed using only demographic and socioeconomic variables.
    \item \textbf{Scenario II:} The analysis was extended to incorporate two variables, defined in a previous section that capture geographical and distance-related factors, including functional economic connectivity as reflected in market potential.
    \end{enumerate}


%{\color{red} EMILI: ?`stage o scenario? No es lo mismo.}




In both clustering processes, the Silhouette {\color{red} (Silhouette en may\'usculas o min\'usculas?)} scores indicated that the optimal number of clusters lies between five and six. To enhance resolution and achieve greater granularity in the classification, we selected a six-cluster solution {\color{purple} \citep{januzaj2023determining}}. {\color{red} (Citas Kolari and Zardkoohi)}

This choice yields a Silhouette coefficient of $0.34$ for the clustering without geographical and distance-related variables (Scenario I), and $0.31$ when these factors are included (Scenario II). While these values do not exceed the conventional threshold of $0.5$, which typically signals a clear separation between clusters, they are nonetheless sufficient to consider the classification robust and useful for the socioeconomic interpretation of the territory.

\end{comment}


\section{Clustering Results}
\label{sec:results}

Following the methodology described in the previous section, our empirical strategy relies on partition-based cluster analysis (k-means) implemented in two scenarios. In the first, clusters are identified using demographic and socioeconomic variables. In the second scenario, geographic variables are incorporated to account for spatial context and functional connectivity. Subsequently, the clustering results obtained with and without geographic variables are compared. The resulting classifications in both cases are evaluated using several validation criteria, including an analysis of municipal reclassifications and quantitative assessments of the robustness and meaningfulness of the LBP classification. Finally, the distribution of municipalities and their populations across the identified clusters is examined.


Regarding the number of clusters, in both scenarious, the results consistently indicate that $K=6$  provides the best balance between within-cluster cohesion and between-cluster separation. The corresponding silhouette and elbow profiles  are reported in Figures \ref{fig:silhouette} and \ref{fig:elbow} of Appendix \ref{sec:appendix_B}.


\subsection{Classification of municipalities}


The map represented in Figure \ref{fig1:clustermap} displays the spatial distribution of clusters based solely on demographic and socioeconomic variables  (scenario I). In contrast, Figure \ref{fig2:clustermap} illustrates the distribution of clusters when all three types of variables (demographic, socioeconomic and geographic) are also included  (scenario II).



To assess the internal consistency of the clustering solutions, we perform a multivariate analysis of variance (MANOVA) for each scenario. This approach allows us to evaluate whether the clusters identified under each specification exhibit distinct multivariate profiles across the set of variables considered. The results, reported in Tables \ref{tab:manova_nogeo} and \ref{tab:manova_geo}, indicate highly significant differences across clusters in both scenarios. In the baseline specification (without geographical variables), Wilks’ Lambda is 0.0027 (F = 2219.3, p < 0.001), while in the extended specification (including geographical variables) it remains similarly low at 0.0033 (F = 1656.2, p < 0.001). These findings confirm that, in both cases, the clustering solutions yield well-defined and statistically distinct groups. While the inclusion of geographic variables modifies the composition of clusters, it does not weaken their internal differentiation, suggesting that the extended specification captures additional dimensions of territorial heterogeneity without compromising the overall structure.



Tables \ref{tbl:2} and \ref{tab:cluster_means}  report the average values of the variables for each cluster in their respective analyses. We begin by examining the first classification (based only on demographic and socioeconomic variables), highlighting the key characteristics of each cluster and identifying their main features and patterns of growth or decline. \footnote{From Table 2, the main characteristics of the clusters are the following: \textit{Cluster 1} (dark red in Figure \ref{fig1:clustermap}) captures metropolitan cores characterized by high income levels and relatively young demographic structures. This is the densest cluster, displaying strong population growth (0.16) and the highest income index (1.04). As shown in Table \ref{tbl:2}, it also exhibits the lowest values for the agricultural sector share (0.03) and the elderly dependency ratio (0.29) among all clusters. The industrial sector accounts for 12\% of employment, a share only lower than that observed in \textit{Cluster 2}. \textit{Cluster 2} (yellow in Figure \ref{fig1:clustermap}) is likely composed of semi-urban or industrial towns characterized by relative economic stability. This group exhibits moderate population density (119.92), near-zero population growth (0.01), and slightly above-average income levels (1.02). One of its most distinctive features is the highest share of industrial employment (0.50), while the agricultural sector accounts for only a small share (0.07). In demographic terms, the share of foreign-born population is moderate (0.14), and the elderly dependency ratio is also at an intermediate level (0.40). \textit{Cluster 3} (purple in Figure \ref{fig1:clustermap}) represents urbanising areas with dynamic demographic patterns, likely corresponding to peri-urban municipalities. This cluster exhibits medium–high population density (160.11) and strong positive population growth (0.23). Income levels are around the average (1.00). In terms of productive structure, both agriculture (0.06) and industry (0.07) account for relatively small shares. Notably, this cluster shows the highest share of foreign youth (0.17) and a relatively low elderly dependency ratio (0.36). \textit{Cluster 4} (green in Figure \ref{fig1:clustermap}) is mainly composed of rural areas characterised by an ageing population and a strong agricultural base. This cluster shows negative population growth relative to the Autonomous Community and a low relative income index. The share of agricultural employment is very high (0.54), while the industrial share remains low (0.05). Migration rates are low, as is the share of foreign youth (0.13), whereas the elderly dependency ratio is relatively high (0.44). \textit{Cluster 5} (blue in Figure \ref{fig1:clustermap}) likely represents small towns in demographic and economic decline, characterised by a strongly ageing population (0.57) and limited dynamism. This group shows low population density (12.66), negative population growth (-0.32), and below-average income levels (0.95). In terms of productive structure, the agricultural sector has a moderate presence (0.21), while industrial employment also represents a relatively small share (0.08). \textit{Cluster 6} (orange in Figure \ref{fig1:clustermap}) is composed of isolated rural areas facing acute demographic decline and economic stagnation. This is the most sparsely populated cluster (7.99/km$^{2}$), with the lowest population growth (-0.34), low relative income levels (1.01), and very limited access to banking services. The shares of both agriculture and industry are very low, while the elderly dependency ratio is the highest (0.74) and the share of foreign-born youth is extremely low (0.03).}






% Summarizing,

% \begin{itemize}
%     \item \textbf{Clusters 1 and 3} represent urban and peri-urban territories, with high density, positive population trends, higher income, and a younger age structure.
%     \item \textbf{Cluster 2} stands out for its strong industrial base and moderate demography, suggesting economic functionality in intermediate areas.
%     \item \textbf{Cluster 4} represents traditional rural zones, with dominant agriculture, ageing, and economic vulnerability.
%     \item \textbf{Clusters 5 and 6} show clear signs of demographic and economic distress, including population loss, low density, high ageing, and limited economic activity.
% \end{itemize}

% \begin{description}
%  \item[Clusters 1 and 3] represent urban and peri-urban territories, with high density, positive population trends, higher income, and a younger age structure.
%
%  \item[Cluster 2] stands out for its strong industrial base and moderate demography, suggesting economic functionality in intermediate areas.
%
%  \item[Cluster 4] represents traditional rural zones, with dominant agriculture, ageing, and economic vulnerability.
%
%  \item[Clusters 5 and 6] show clear signs of demographic and economic distress, including population loss, low density, high ageing, and limited economic activity.
% \end{description}
%

Therefore, it could be possible to summarise this typology as: (i) clusters 1 and 3 represent urban and peri-urban territories, with high density, positive population trends, higher income, and a younger age structure; (ii) cluster 2 stands out for its strong industrial base and moderate demography, suggesting economic functionality in intermediate areas; (iii) cluster 4 represents traditional rural zones, with dominant agriculture, ageing, and economic vulnerability; and (iv) clusters 5 and 6 show clear signs of demographic and economic distress, including population loss, low density, high ageing, and limited economic activity.



In this context, Clusters 5 and 6 could be emerging as the most representative examples of left-behind places. Cluster 6 is particularly critical, characterised by extremely low population density, population decline, a high dependency ratio, a weak economic base, and an almost complete absence of services (e.g., banking infrastructure). Cluster 5 also shows clear signs of vulnerability, including stagnant population growth, low income levels, high ageing rates, and a limited industrial and agricultural base, although it benefits from slightly better access to services than Cluster 6. Cluster 4, being composed of LBP municipalities, could also be considered a borderline case. \footnote{Depending on whether its strong agricultural orientation is interpreted as a source of resilience \citep{martin} or, conversely, as a structural constraint.}


In the second scenario we apply the same clustering methodology but include geographic variables. Specifically, we also include proximity to urban or semi-urban centres (as measured by equation \eqref{eq:eq9}) and a measure of potential market income \citep{harris1954market} and Krugman (1991). This allows for a more detailed understanding of territorial disparities. Figure \ref{fig2:clustermap} presents the resulting cluster map with the inclusion of these geographic variables. Table \ref{tab:cluster_means} shows the average values of the variables for each new cluster.\footnote{These clusters are described as follows:
\textit{Cluster 1} (dark red in Figure \ref{fig2:clustermap}) represents the most urbanized municipalities. These areas display by far the highest population density (2563.87 inhabitants/km$^{2}$), the shortest distance to urban and semi-urban centres (0.11), and the highest market potential (139,243,952). Demographically, they show a relatively low elderly dependency ratio (0.29) together with moderate positive population growth (0.11). From an economic perspective, agriculture plays only a marginal role (0.03), while industry has a moderate presence (0.13). Overall, these characteristics clearly identify this cluster as metropolitan cores with strong accessibility and economic integration. \textit{Cluster 2} (yellow in Figure \ref{fig2:clustermap}) groups municipalities characterised by a strong industrial base and good accessibility. The industrial employment share is the highest among all clusters (0.50), combined with above-average income levels (1.02). Population density (118.42 inhabitants/km$^{2}$) and accessibility indicators place these municipalities in intermediate but well-connected locations, with a moderate distance to urban or semiurban centres (6.36). Demographic indicators are relatively balanced, with moderate ageing (0.39) and slightly positive migration dynamics. Overall, this cluster represents economically dynamic municipalities with good connectivity and a clear industrial profile. \textit{Cluster 3} (purple in Figure \ref{fig2:clustermap}) includes municipalities with the strongest demographic dynamism, showing the highest population growth (0.46). These municipalities combine relatively high density (206.67 inhabitants/km$^{2}$) with good accessibility to urban centres (distance index 4.57) and high market potential (117,940,928). Income levels are also above the national average (1.08). Sectorally, the economy is relatively diversified, with moderate shares of both industry (0.08) and agriculture (0.07). Altogether, this cluster reflects growing municipalities located in relatively central and economically favourable environments. \textit{Cluster 4} (green in Figure \ref{fig2:clustermap}) corresponds to traditional rural municipalities with a strong agricultural specialization. This cluster records the highest agricultural employment share (0.36), while population density remains moderate (66.98 inhabitants/km$^{2}$). Accessibility is more limited than in the previous clusters, with a distance index of 9.04 and comparatively low market potential (46,838,040). Demographically, these municipalities show signs of ageing (dependency ratio 0.43) and negative population growth (-0.15). These characteristics describe traditional rural areas with moderate accessibility and a clear agricultural economic structure. \textit{Cluster 5} (blue in Figure \ref{fig2:clustermap}) is composed of structurally vulnerable rural municipalities. These territories present low population density (12.32 inhabitants/km$^{2}$), an ageing demographic structure (dependency ratio 0.56), and clear population decline (-0.31). Income levels are below the national average (0.94) and accessibility constraints are significant, as reflected in the highest distance to urban centres (17.73). The agricultural sector remains relatively relevant (0.22), although less dominant than in Cluster 4. Despite displaying moderate market potential (62,413,868), their remoteness likely limits effective access to economic opportunities, reinforcing their condition as vulnerable rural areas. Finally, \textit{Cluster 6} (orange in Figure \ref{fig2:clustermap}) represents the most critical rural periphery. These municipalities exhibit extremely low population density (6.55 inhabitants/km$^{2}$), the highest elderly dependency ratio (0.75), and strong population decline (-0.34). Economic activity is weak and poorly diversified, with very low shares of both industry (0.02) and agriculture (0.06), together with an almost complete absence of banking services. Accessibility constraints are also severe, with a very high distance to urban and semi-urban centres (17.51). Although market potential is not the lowest (88,904,232), this likely reflects proximity to more prosperous areas without effective accessibility. Overall, this cluster represents the most acute manifestation of left-behind rural territories.

}



In Table \ref{tbl:5}, the main cluster characteristics and their corresponding vulnerability status, based on the results reported in Table \ref{tab:cluster_means}, are summarized. As observed, CL1, CL2 and CL3 correspond to urban, industrial, accessible and economically dynamic municipalities, displaying relatively favourable demographic and socioeconomic conditions. By contrast, CL4, CL5 and CL6 can be classified as left-behind places, although with differing degrees and forms of vulnerability. Specifically, CL4 comprises rural and agricultural territories with moderate accessibility and emerging signs of vulnerability; CL5 reflects structurally vulnerable rural areas characterized by demographic decline and weak integration into major economic centres; while CL6 represents the most critical situations, marked by extreme territorial isolation, severe demographic deterioration and a very limited provision of services.





\subsection{Comparison of clustering results with and without geographic variables}



The comparison between the two clustering scenarios underscores the importance of geography in the definition of left-behind places. Indicators such as distance to semi-urban and urban centres and measures of market potential capture differences in functional connectivity between municipalities, shedding light on how spatial isolation or proximity to economic cores may condition local development trajectories.\citep{connor2024gets}. When the geographic dimension is excluded, key inter-territorial linkages may remain hidden, and part of the underlying heterogeneity may be obscured if only demographic and socio-economic indicators are considered.

Beyond confirming the relevance of geography, this two scenarios clustering strategy provides additional insights into the contribution of second-nature geographic variables to the identification of left-behind places. In particular, it allows us to examine the extent to which municipalities change their classification once second-nature geography is taken into account. This study places special emphasis on municipalities that transition from a non-LBP to an LBP classification. These transitions reveal cases where demographic and socioeconomic characteristics alone mask underlying vulnerabilities that become visible only when geographical constraints and functional disconnection are incorporated into the analysis.

Against this background, both the map in Figure \ref{fig3:clustermap} and Table \ref{tab:reassign_clusters} illustrate how the cluster classification of municipalities changes when geographic variables are included in the analysis. From Figure \ref{fig3:clustermap} and Table \ref{tab:reassign_clusters}, municipalities shown in white (6,496 municipalities) remain in the same cluster, indicating that the incorporation of spatial criteria does not alter their classification. However, due to the inclusion of these new variables,  1,601 municipalities (around 20\%) are reassigned to a different cluster. Of the 1,601 reassigned municipalities, 75.7\% (1,211 municipalities) change their status regarding whether they are classified as LBP or not. Among the 1,211 municipalities that change their LBP status, 11.7\% (142 municipalities) shift from LBP to non-LBP, whereas the remaining 88.3\% (1,069 municipalities) move from non-LBP to LBP \footnote{A significant number of these (1,044 municipalities) were initially classified in Cluster 3, which includes urban and peri-urban areas (Figure \ref{fig1:clustermap}).
Of these, 751 are reassigned to Cluster 4 (Traditional rural/agricultural with moderate accessibility), 26 to cluster 5 (structural and geographical vulnerability) and 233 to Cluster 6 (critical rural periphery). The remaining 34 municipalities are reclassified into clusters other than Cluster~3 but do not transition into left-behind positions. Reclassification effects are also observed among municipalities initially assigned to Cluster~2 (industrial, accessible areas with solid income levels). In this case, 35 municipalities are reassigned to Cluster~4 (rural and agricultural, with moderate accesibility), 8 to Cluster~5 (structural and geographical vulnerability), and 13 to Cluster~6 (extreme isolation, strong demographic decline and limited service provision). The remaining 34 municipalities are reclassified into non–left-behind clusters distinct from Cluster~2.
}.





Therefore, the introduction of second-nature geographic variables reveals that many of these municipalities (142 of 1601, that is, 8.9\%) are less favourably located than their socio-economic indicators alone would suggest. Despite displaying dynamic characteristics, their weak territorial integration and limited inter-municipal connectivity undermine their development potential and lead to a downgrade in their cluster assignment \citep{torre2014regional}.\footnote{For example, this the case of Badajoz, a Spanish capital city and the most important urban center in the region of Extremadura. When clustering only with demographic and socioeconomic variables, Badajoz is classified in Cluster 3, characterized by urban and peri-urban areas with high population density, favorable demographic trends, higher income levels, and younger age structures. Nevertheless, when geographic variables are included, Badajoz is reassigned to Cluster 4, corresponding to rural/agricultural municipalities with moderate accessibility. Badajoz shifts from a cluster representing more urban characteristics to one reflecting rural and agricultural features. This would imply that a correct characterization of a municipality would require the consideration of geographic variables.} %However, when geographic variables such as distance to urban centers (dist_1_std = -1.31)
%and market potential (PM_renda_std = -1.18) are included. Besides the low market potential, other variables contribute to this reclassification. Notably, the industrial employment share is relatively low (share\_ind\_std = -0.36), indicating a less industrialized local economy typical of rural areas. Additionally, the dependency ratio for elderly population is also significantly below average (dep\_mayor\_std = -0.93), and population density is moderately lower than average (dens\_pop\_std = -0.08). These factors combined with the geographic variables explain why Badajoz shifts from a cluster representing more urban characteristics to one reflecting rural and agricultural features. {\color{red} Again, this would imply...}
 On the other hand, around 66.8\% of municipalities (1,069 municipalities) change their cluster assignment moving into a left-behind position once spatial factors are taken into account. The remaining 24.4\% of municipalities are reclassified without altering their left-behind status, either remaining within left-behind or non–left-behind categories. In sum, the inclusion of second-nature geographic variables not only shifts the clustering structure but proves essential for a more accurate identification of territorially disadvantaged areas. These findings reinforce the importance of spatial positioning and interterritorial linkages in local development.

In addition to the qualitative description of the variations observed in the classification of municipalities when geographic and spatial accessibility variables are included, quantitative indicators underscore the significant influence of these variables on the clustering structure.



\begin{comment}
Other indicator is the Rand Index {\color{red} (¿se ha hablado de él en algún momento? Queda un poco pegote)}, which measures the similarity between two different partitions of a dataset \footnote{This index ranges from 0 to 1, where 0 indicates completely dissimilar partitions and 1 indicates identical ones. In our analysis, the Rand Index is 0.56,}, suggest a moderate level of similarity between the classification based solely on demographic and socioeconomic variables and the classification that incorporates geographic characteristics. This result confirms that the inclusion of spatial features introduces a new dimension that substantially alters the categorization, enabling a more nuanced identification of municipalities that can be considered ``left behind places'' \citep{rand1971objective}.




Regarding the Silhouette Scores, the clustering without geographic variables yields 0.34, while the model including geographic and spatial accessibility variables scores 0.31. This difference is small and unlikely to be statistically significant, but it highlights changes in the structure and interpretation of the clusters. Both values indicate a moderate clustering with reasonably well-separated groups. The slightly lower score in the model with geographic variables suggests more complex or overlapping cluster boundaries; however, the additional explanatory variables contribute to more coherent and interpretable profiles.



Furthermore, the Calinski-Harabasz Index \citep{calinski1974dendrite} supports this conclusion, with a value of 1100.81 for the clustering without geographic variables and a higher value of 1164.89 when geographic and spatial accessibility variables are included. The increase in this index indicates improved cluster cohesion and separation, reinforcing the idea that the integration of geographic variables contributes to a more robust and meaningful classification of municipalities .

\end{comment}


To assess the robustness of the clustering results, the classifications obtained under both scenarios are compared using the following indicators. First, the average silhouette value of the final clustering solution is used to assess overall clustering quality. In our case, this value equals 0.34 for the specification excluding geographical variables and 0.32 for the specification including geographical variables, indicating a reasonably well-defined clustering structure in both cases.\footnote{Here, the average silhouette value refers to the mean silhouette coefficient computed across all municipalities for the final partition with $K=6$. Appendix \ref{sec:appendix_B} reports a detailed description of silhouette index.}
Second, the similarity between the two clustering specifications is evaluated using the Adjusted Rand Index (ARI), which measures the extent to which municipalities are assigned to the same clusters across the two classifications. The resulting value (ARI = 0.56) indicates a moderate level of agreement, suggesting that while the overall cluster structure remains relatively stable, the inclusion of geographic variables leads to a partial reallocation of municipalities across clusters \footnote{The Adjusted Rand Index (ARI) measures the similarity between two classifications by comparing the agreement in the assignment of observations to clusters while correcting for similarity that may occur by chance. The index ranges from $-1$ to $1$, where values close to $1$ indicate a high level of agreement between classifications, values around $0$ suggest random similarity, and negative values indicate systematic disagreement.}. Finally, spatial coherence is examined using Moran's $I$ statistic, which measures the degree of spatial autocorrelation in cluster assignments and therefore indicates whether municipalities belonging to the same cluster tend to be geographically contiguous. It is noteworthy that first-order spatial autocorrelation emerges when Moran’s I is computed for the cluster assignment obtained without including geographic variables. Table \ref{tab:moran} shows that Moran's $I$ increases from 0.334 in the specification based solely on socio-economic variables (with a high level of statistical significance) to 0.419 when geographical variables are incorporated, revealing a stronger spatial coherence. The inclusion of geographical variables such as distance to urban centres and market potential does not introduce a spurious spatial dependence in the clustering procedure.\footnote{Rather than incorporating geographical coordinates or spatial weight matrices, these indicators capture distance-related frictions and economic accessibility conditions associated with second nature geography.} The slight increase in Moran’s $I$ observed when these variables are included suggests that the resulting classification becomes more spatially coherent, without substantially altering the overall cluster structure. This indicates that accessibility variables complement, rather than dominate, the socio-economic and demographic information used in the clustering procedure.
Thus, our geographic variables do not introduce spatial autocorrelation; they reveal the spatial-economic structure already embedded in the territorial system.

Taken together, the qualitative and quantitative evidence suggests that the classification incorporating geographic variables should be preferred. This approach is consistent with recent literature on left-behind regions, which emphasises the need to incorporate greater “geographical wisdom” when identifying regional typologies {\color{red} (Jessen, 2024)}. By accounting for spatial accessibility conditions and territorial context, the classification better reflects the spatial-economic mechanisms underlying regional divergence. Accordingly, the remainder of the paper focuses on the results obtained under Scenario 2, which includes the geographic variables described above.








\subsection{Regional distribution and composition of clusters: Where people live vs where municipalities are}

In this section, the analysis is aggregated from the municipal level to the level of Autonomous Communities, which correspond to the NUTS-2 classification used by the European statistical system. While municipalities provide a highly detailed spatial perspective, the NUTS-2 level represents the standard territorial scale for regional analysis and policy design within the European Union. Aggregating the results at this level therefore enables the findings to be interpreted within a broader regional context and facilitates their potential use in regional development and cohesion policy frameworks.



Table \ref{tbl:7} reports the regional composition of each cluster in terms of the share of municipalities and population (column percentages). The results reveal a marked asymmetry between the territorial and demographic distribution of clusters. While clusters CL4--CL6 (LBPs) account for over 71\% of Spanish municipalities, they represent only about 31\% of the total population. By contrast, clusters CL1--CL3 (NLBPs) concentrate nearly 69\% of the population despite comprising less than 29\% of municipalities. The table also shows a clear regional concentration of several clusters, with some regions contributing disproportionately to specific cluster types. In particular, more than half of the regional population resides in LBP municipalities in Andalusia, the Canary Islands, Castile--La Mancha, Extremadura, and the Region of Murcia. By contrast, highly urbanised and metropolitan regions such as the Community of Madrid, Catalonia, and the Basque Country are predominantly composed of non-LBP municipalities.\footnote{NLBP clusters are dominated by the large metropolitan areas: CL1 (urban, hyper-connected) population is led by Madrid and Catalonia (each $\approx 26-27\%$), followed by the Valencian Community and Andalusia; CL2 (industrial, accessible) concentrates along the Mediterranean corridor and the northern belt (Catalonia, Valencian Community, Basque Country, Galicia, Aragón); CL3 (dynamic, central) is more distributed across Catalonia, Madrid, Valencian Community, Basque Country and Castile and León. LBP clusters concentrate outside the big metros: CL4 (rural/agricultural, moderate accessibility) is dominated by Andalusia ($\approx 37\%$ of its population) with sizable shares from the Valencian Community, Murcia, Canary Islands, Castilla--La Mancha, Galicia and Extremadura; CL5 (structural/geographical vulnerability) concentrates in Castile and León and Aragón; CL6 (critical isolation, no services) is centered in Castile and León, Galicia and Asturias (with contributions from Castilla--La Mancha and Catalonia). Nationally, around \(71\%\) of municipalities are LBP while \(\sim 69\%\) of people live in NLBP—i.e., LBP is \textsl{municipality-intensive} whereas NLBP is \textsl  {population-intensive}.}

Table \ref{tbl:6} shows the internal composition of each of the 17 Spanish regions in terms of LBP and non-LBP municipalities, based on the cluster classification of their municipalities. Deriving this regional pattern through a bottom-up approach—by aggregating municipalities according to their cluster classification—constitutes a novel contribution to the literature. The results reveal substantial regional heterogeneity in the prevalence of LBPs. In several regions—such as Extremadura, Castile--La Mancha, Andalusia, the Canary Islands, and the Region of Murcia—the majority of municipalities and population are located in LBP clusters. By contrast, metropolitan and economically dynamic regions such as the Community of Madrid, the Basque Country, and Catalonia are overwhelmingly composed of non-LBP municipalities, both in terms of population and municipal counts. These patterns highlight the strong territorial unevenness underlying the spatial distribution of left-behind places in Spain.\footnote{Table \ref{tbl:6} reports the distribution of municipalities and population across clusters by Autonomous Community, reflecting the regional context of Spain. Table \ref{tbl:6} shows row shares (100\%~= CCAA) and reveals a clear split between metropolitan regions and rural interiors. Southern and interior regions concentrate most of their \textsl{municipalities} in LBP clusters---e.g., Extremadura ($\approx 98\%$) of municipalities, $\approx 90\%$) of population in LBP, mainly CL4),
Castilla--La Mancha, Andalusia, Canary Islands, Murcia, and Galicia. By contrast, metropolitan and industrial regions place the bulk of their \textsl{population} in NLBP clusters---Madrid ($\approx 99\%$), Catalonia, Basque Country and Navarre.
An important asymmetry appears in Castile and León: $\approx 84\%$ of its municipalities fall in LBP (especially CL6, \textsl{extreme isolation/decline}), yet $\approx 68\%$ of its population remains in NLBP.}




Together, Tables \ref{tbl:7} and \ref{tbl:6} map a dual geography: a metropolitan/industrial system (CL1--CL3) where most people live, and a rural/vulnerable system (CL4--CL6) where most municipalities lie, with CL6 pockets requiring priority attention despite small populations. Moreover, for regions in which most municipalities and population belong to LBP clusters, the results uncover a particularly severe warning signal—largely overlooked in previous research—that should be taken into account in the formulation of economic policy. \footnote{Although, to the best of our knowledge, this is the first study to estimate the proportion of Spanish municipalities and population that can be classified as left-behind places (LBPs), previous contributions such as  \cite{bde2021distribucion}, \cite{villamayor2024espana}, among others, do not fully reveal the magnitude of the territorial challenges identified in our analysis.}

Figure \ref{fig:lbp_scatter} plots the relationship between the share of municipalities classified as LBPs and the share of population living in LBP municipalities across Spanish regions. The figure reveals substantial regional heterogeneity and a generally positive association between the two variables: regions with a higher proportion of LBP municipalities also tend to concentrate a larger share of their population in LBP areas. However, important differences emerge. Some regions, such as Extremadura, the Region of Murcia, and Castile--La Mancha, combine very high shares of LBP municipalities with a similarly high share of population living in such areas. By contrast, metropolitan and economically dynamic regions such as the Community of Madrid, the Basque Country, and Catalonia display very low shares of population residing in LBP municipalities despite having some municipalities classified as LBPs. These patterns further illustrate the strong regional unevenness in the spatial distribution of left-behind places in Spain.


Table \ref{tbl:8} presents the distribution of municipalities and population by degree of urbanisation within each cluster. The results highlight clear structural differences. Urban municipalities are mostly concentrated in Cluster 1 (61.35\% of municipalities and 71.77\% of the population), while smaller shares are distributed across Clusters 2--4. Intermediate municipalities are mainly found in Clusters 3 and 4 (36.41\% and 41.94\% of municipalities, respectively), which together account for over 80\% of the intermediate population. Rural municipalities are spread across almost all clusters, though Clusters 4 and 6 concentrate most of them (31.02\% and 38.80\%, respectively). Population-wise, Cluster 4 stands out with more than half (51.99\%) of the rural population.
%\end{itemize}



Table \ref{tbl:8} confirms that urban municipalities, although fewer in number (215 out of 8,097), account for the largest share of the population (19.7 million out of 47.2 million), while rural municipalities dominate in numerical terms (6,761 municipalities) but only host about 6.2 million inhabitants.

Table \ref{tbl:9} shows a complementary perspective: the distribution of municipalities and population across clusters within each degree of urbanisation. Among urban municipalities, most belong to Cluster 1 (71.63\%), reflecting its strong demographic weight, while Clusters 3 and 4 together host nearly a quarter of the urban municipalities. Intermediate municipalities are more evenly distributed, with Clusters 2, 3, and 4 each gathering significant shares (24--28\%). Finally, rural municipalities are highly concentrated in Clusters 2, 3, 4, and 6, with Cluster 6 alone accounting for over one third of rural municipalities; in population terms, Cluster 4 is again the most relevant, containing almost 24\% of rural residents.




Overall, the combination of Tables \ref{tbl:8} and \ref{tbl:9} illustrate the asymmetric relationship between urbanisation and clustering: a small number of urban municipalities (largely in Cluster 1) dominate the demographic structure, while rural municipalities are numerous but fragmented across several clusters with much smaller population shares.





\section{Post-cluster analysis}
\label{sec:post_cluster}
This section presents a post-cluster analysis examining the extent to which demographic, socioeconomic, and geographic dimensions influence cluster formation. We begin with an analysis of boundary sensitivity (Subsection \ref{sec:anova}), based on thresholds and decision boundaries, which examines how close municipalities are to transitioning from one cluster to another and which dimensions operate as the most immediate constraints on such changes. This approach focuses on the geometry and stability of the classification, making it possible to evaluate the robustness of the clustering structure and to identify the factors that are most decisive at the margins between groups. We then turn to the multinomial logit model Subsection \ref{sec:multinomial}, which explores how different variables are associated with the probability of belonging to each cluster relative to a reference category. Whereas the boundary sensitivity analysis captures proximity to cluster transitions and the stability of classification, the multinomial model provides a probabilistic and associative interpretation of cluster membership. Taken together, both approaches offer complementary insights into the explanatory mechanisms and robustness of the clustering results.

Taken together, both approaches offer complementary and mutually reinforcing insights: the frontier-based analysis identifies the dimensions along which municipalities are closest to changing cluster, whereas the multinomial model highlights the variables that most strongly influence cluster membership. Their joint consideration allows us to uncover the structural mechanisms underlying each type of left-behind place and to assess the robustness of the clustering results.







\subsection{Robustness and sensitivity analysis:  thresholds and decision boundaries}
\label{sec:anova}



As a first step, it was verified that the resulting clusters are not the result of random variation \citep{everitt2011cluster}. To this end, a one-way ANOVA was performed for each variable across the six clusters. The results, reported in Table \ref{tbl:anova}, indicate that all variables exhibit statistically significant differences across clusters (p < 0.001), confirming the robustness of the clustering structure.\footnote{Levene’s tests reject the assumption of homogeneity of variances for all variables (p < 0.001), indicating the presence of heteroscedasticity across clusters. This implies that clusters differ not only in their mean characteristics but also in their internal variability. For this reason, Welch’s ANOVA is used as a robustness check, as it does not rely on the assumption of equal variances. The results remain qualitatively unchanged, confirming that the observed differences are not driven by variance heterogeneity.}








To further assess the role of individual dimensions in shaping cluster membership, the study proposes a novel empirical application of post-clustering sensitivity analysis based on threshold calculations. \citep{hastie2009elements}, \citep{bishop2006pattern}.\footnote{  This exercise is designed to examine the geometric proximity of municipalities to cluster boundaries in the multidimensional space defined by the clustering procedure. Its purpose is not to identify causal mechanisms or policy levers, but to assess which dimensions are most closely associated with marginal changes in cluster membership.

Specifically, thresholds indicate the value that a given variable would need to reach, holding all other characteristics constant, for a municipality to lie exactly on the decision boundary between its current cluster and the closest alternative cluster. This framework allows us to evaluate whether particular dimensions mechanically dominate the geometry of the clustering space, a concern frequently raised in relation to the inclusion of geographical variables.

The analysis focuses on Scenario II (including geographic variables), and is implemented asymmetrically, in line with the conceptual distinction between left-behind and non-left-behind places. For municipalities classified as left-behind (Clusters~4, 5 and~6), thresholds are computed with respect to the closest non-left-behind cluster (Clusters~1, 2 or~3), thereby capturing proximity to exiting left-behind status. Conversely, for non-left-behind municipalities, thresholds are defined relative to the closest left-behind cluster. This approach ensures that the sensitivity analysis focuses on transitions that are substantively meaningful for the research question, rather than on reallocations within the same broad category.

Thresholds are calculated for each municipality and each variable included in the clustering procedure. Two related measures are derived. First, the threshold value identifies the critical level that the variable would need to attain for the municipality to lie exactly on the decision boundary between its current cluster and the alternative cluster. Second, the corresponding distance-to-threshold measure captures the magnitude of change required relative to the municipality’s observed value. As all variables are standardised prior to clustering, these distances can be interpreted as relative displacements in a common metric, allowing direct comparison across dimensions.}

Focusing on Scenario II (including economic, demographic and geographic dimensions), we next analyse which dimension(s) most frequently drive changes in cluster membership. Therefore, Table~\ref{tab:threshold_dimension} reports the dimension associated with the closest decision boundary for municipalities classified as left-behind and for those that switch to left-behind status when geographical variables are introduced. Two results stand out. First, the geographic variables never define the closest decision boundary, either for left-behind municipalities as a whole or for those that change classification. This indicates that distance and market potential do not mechanically dominate the clustering geometry once demographic and socioeconomic characteristics are taken into account. In other words, geographic variables do not act as binding constraints at the margin of cluster membership. Second, among municipalities that switch to left-behind status when geographic variables are incorporated, the closest decision boundaries are overwhelmingly associated with socioeconomic dimensions. This pattern suggests that the inclusion of geography primarily reveals latent structural vulnerabilities
rather than artificially classifying municipalities as disadvantaged simply because of distance.\footnote{When analysing which dimension most frequently drives changes in cluster membership, the results indicate that the economic dimension is the most influential, followed by the demography. This does not imply that the remaining dimensions are not considered in the analysis, but rather that they rarely act as the binding constraint. This finding reinforces an important point: although incorporating geography increases the number of municipalities classified as LBP, the immediate trigger of reclassification is typically not an accessibility variable, but primarily income (and, to a lesser extent, population growth). Overall, this suggests that municipalities transitioning into LBP status—as a result of including geographical factors—were already exposed to less favourable spatial conditions, but the change in classification is mainly activated through economic (income) and demographic (population growth) factors.}

Importantly, the dimension that defines the closest decision boundary should not be interpreted as the main structural constraint underlying left-behind status. Rather, it identifies the margin along which a municipality is least distant from an alternative cluster, given its overall configuration of characteristics. When income- or other socioeconomic-related variables define the closest boundary, this indicates that such dimensions operate as marginal adjustment margins, while deeper structural constraints are likely to be associated with dimensions characterised by much larger threshold distances.

%Figure~\ref{fig:dthr_dimension} displays the distribution of absolute distances to the closest decision boundary for the dimensions that effectively define marginal transitions. Accessibility-related variables do not appear in the figure because they never constitute the closest decision boundary for left-behind municipalities, a result consistent with the evidence reported in Table~\ref{tab:threshold_dimension}.

Overall, this post-clustering sensitivity analysis reinforces the interpretation that the inclusion of geographical and accessibility-related variables refines, rather than distorts, the identification of left-behind places. While geography plays a role in shaping cluster membership, it rarely determines marginal transitions between clusters. Instead, left-behind status emerges as the outcome of multidimensional disadvantage, in which functional accessibility interacts with demographic and socioeconomic factors without mechanically driving the resulting typology.

A natural question arising from this analysis is why a clustering scenario that includes second-nature geographic variables is necessary if the geographic dimension rarely defines the closest decision boundaries. The answer lies in the distinction between identification and marginal proximity. The second clustering scenario is designed to identify left-behind places by incorporating second-nature geographic variables into the definition of territorial disadvantage. By contrast, the threshold analysis examines which individual dimension would need to change marginally for a municipality to exit left-behind status. The limited role of second-nature geographic variables in defining the closest boundaries therefore reflects their position in marginal adjustments rather than their contribution to the identification of left-behind places.






To further characterise the distinct types of left-behind places identified in the clustering analysis, Table~\ref{tab:structural_blocks} summarises the median absolute distance to the closest decision boundary (\(|dthr|\)) by dimension for each left-behind cluster \footnote{Table~\ref{tab:threshold_by_cluster} reports the median distance to threshold for each individual variable by LBP cluster, allowing us to identify the specific indicators that lie closest to the classification frontier.}. Larger values indicate dimensions along which municipalities are structurally distant from any marginal transition to alternative clusters. The patterns are strongly heterogeneous across left-behind types. Cluster~4 exhibits its largest structural distances along demographic dimensions, whereas its geographic distances are comparatively smaller, suggesting a predominantly demographic form of left-behindness \footnote{For Cluster~4, the median absolute distance-to-threshold along demographic dimensions reaches 12.7 standardised units, indicating that municipalities in this group are structurally distant from any marginal transition out of left-behind status through demographic change alone. The demographic distance reported in Table~\ref{tab:structural_blocks} does not refer to a single variable, but to the typical magnitude of univariate changes required across demographic indicators for municipalities to reach a decision boundary with the closest non-left-behind cluster. By summarising variable-specific thresholds using the median, the indicator captures how structurally distant municipalities are from any marginal demographic transition, without relying on extreme values or implying that a single demographic intervention would be sufficient.
}. Cluster~5 displays the highest demographic distances of all three clusters, combined with sizeable socioeconomic distances, pointing to a more entrenched and multidimensional disadvantage. In contrast, Cluster~6 is characterised by very large socioeconomic distances and comparatively smaller demographic distances, indicating a type of left-behindness primarily rooted in economic and financial constraints rather than demographic structure. Overall, these results reinforce the interpretation that left-behindness in Spain is not a single condition but a set of structurally distinct profiles.



From a policy perspective, these findings suggest that geographic disadvantage does not operate primarily as an autonomous driver of left-behind dynamics. Rather, second-nature geography appears to reveal or amplify underlying socioeconomic weaknesses. Consequently, policies focused exclusively on improving both distances to urban or semiurban centers and the potential market may have limited effects unless they are accompanied by interventions aimed at strengthening the local socioeconomic structure (including labour market opportunities, productive diversification and human capital formation).




%\subsection{Binary Logit Analysis. Robustness Analysis: The Role of Geographic Variables}

%To further explore the determinants of cluster membership and assess the robustness of our classification, we performed a binary logistic regression analysis. For this purpose, the original clusters obtained using geographic and distance variables were recoded into a binary dependent variable, where clusters 1, 5, and 6 were combined to represent Left Behind Places (LBPs) and coded as 1, while all other clusters were grouped together and coded as 0. The logistic regression was conducted using the second clustering solution Map \ref{fig2:clustermap}, which incorporates geographic and spatial accessibility variables, as it was deemed the most appropriate classification \citep{katos2007network}.

%This binary classification allows us to focus on distinguishing LBPs from other municipalities and to assess the individual contribution and statistical significance of each explanatory variable in predicting the likelihood that a municipality belongs to the group of LBPs.

%Table \ref{tbl:8} presents the average marginal effects (AMEs) estimated from the binary logistic regression model based on the classification that includes geographic and accessibility variables (Map 2). These effects quantify the change in the predicted probability of a municipality being classified as a Left Behind Place (LBP) associated with a one-unit increase in each explanatory variable, holding all other variables constant. The analysis highlights which variables most strongly influence the likelihood of being identified as an LBP, providing further insights into the territorial dynamics underlying this classification.

%The marginal effects reported in indicate the average change in the probability of being classified as a left-behind place (LBP) associated with a one-unit increase in each predictor, holding other variables constant. Since the predictors are expressed in proportions (from 0 to 1), a one-unit increase corresponds to a 100 percentage-point increase.

%Among the variables negatively associated with the probability of being classified as an LBP, the most influential are the industrial employment share (share\_ind), the net migration rate (t\_migr), the income per capita of municipality in relation to the country income, youth out-migration share (shext16) and the the relative population growth rate (regpopgrowccaa). For example, a 10 percentage point increase in youth out-migration is associated with a 2.8 percentage point decrease in the probability that a municipality is classified as a Left Behind Place, highlighting the protective effect of a stronger demographic renewal dynamic. Similarly,  a 10 percentage point increase in the share of industrial employment is associated with a 5.9 percentage point decrease in the probability that a municipality is classified as a Left Behind Place. This suggests that a stronger industrial base significantly reduces the likelihood of marginalisation.These negative signs are consistent with the theoretical expectation that stronger demographic dynamics (population growth, in-migration) and greater industrial development reduce the likelihood of socioeconomic marginalisation.

%On the other hand, the agricultural employment share (share\_agr) and the elderly dependency ratio (dep\_mayor) show strong positive marginal effects. A 10-percentage-point increase in the agricultural employment share, for example, raises the probability of being considered a Left Behind Place by approximately 2.1 percentage points. This positive association suggests that a higher reliance on agricultural employment is linked to greater vulnerability to marginalisation. This means that a 1\% increase in the agricultural employment share is associated with an increase of approximately 0.21 percentage points in the probability of being classified as a left-behind place. These results align with the theoretical expectation that a reliance on agriculture and an ageing population are structural vulnerabilities associated with peripheralisation processes.

%The presence of bank branches per capita (shbancos\_pop), although measured on a different scale, also has a significant and positive effect. When rescaled for interpretation, the results suggest that adding one bank office per 1,000 inhabitants increases the probability of LBP classification by approximately 0.050 percentage points. This positive association could reflect a compensatory dynamic in which financial services are more present in structurally disadvantaged areas precisely due to their greater dependency on face-to-face banking, often linked to demographic ageing and lower levels of digitalisation.

%Geographic and distance-related variables also play a significant role in explaining LBP classification. The distance to the nearest urban centre (dist) and the market potential index (PM\_renda) both exhibit statistically significant effects in the expected direction. Greater distance from urban centres increases the likelihood of being classified as an LBP, while higher market potential—reflecting proximity to wealth and economic opportunities—reduces this probability. These results highlight the spatial disadvantage faced by peripheral municipalities and underscore the importance of accessibility and territorial embeddedness in shaping patterns of territorial exclusion.

%Overall, these results confirm the critical role of structural economic, demographic, and geographic factors in shaping territorial disparities. A higher share of agricultural employment and an ageing population significantly increase the likelihood of being classified as a left-behind place, whereas greater industrial specialisation, population growth, and youth retention reduce this probability. In addition, geographic isolation and poor accessibility—captured by variables such as distance to the nearest urban or semiurban centre and lower market potential—are strongly and positively associated with LBP classification, highlighting the importance of spatial proximity and connectivity in territorial cohesion. The signs and statistical significance of the marginal effects are broadly consistent with expectations derived from regional development theory, reinforcing the view that long-term structural and spatial conditions underpin persistent inequalities across municipalities.

\subsection{Multinomial Logit Analysis. Robustness Analysis: The Role of Geographic Variables}
\label{sec:multinomial}



To better understand the factors shaping the classification obtained through K-Means clustering, we estimated a multinomial logistic regression model. The dependent variable is a recoded categorical indicator based on the original cluster assignments. Specifically, the three clusters identified as Left Behind Places (4, 5 and 6) were retained as categories 1, 2 and 3, while the remaining three clusters, corresponding to municipalities not classified as LBPs (1, 2 and 3), were grouped together and coded as 0, serving as the reference category. The reference category corresponds to superstar regions (Kemeny and Storper, 2020a), understood as territories that exhibit strong economic growth, generally featuring higher-than-average population density and a concentration of relatively young, highly skilled, and high-income individuals. The explanatory variables used in the regression are the same demographic, economic, and geographic indicators employed in the clustering process \citep{mirzahossein2023clustering}.

The primary objective of this model is not prediction, but to assess the relative influence of each variable in explaining the structural conditions associated with different types of territorial disadvantage. In particular, it allows us to determine whether each LBP profile is primarily associated with demographic, economic, or geographic factors, controlling for the simultaneous effect of all other factors.

We present the results of the regression in the form of odds ratios.\footnote{which provide a more intuitive and policy-relevant measure of effect size than raw coefficients. The odds ratio expresses the multiplicative change in the relative probability (or odds) of belonging to a given LBP category versus the reference group, for a one-unit increase in the explanatory variable:

\begin{equation}
\label{eq:odds_ratio}
\text{Odds Ratio} = e^{\beta}
\end{equation}


\noindent An odds ratio of 1 indicates no effect relative to the non-LBP reference category. Values greater than 1 suggest a higher likelihood of belonging to that LBP group, and values below 1 indicate a lower likelihood.

To ensure comparability and interpretability across variables with different scales, all explanatory variables were standardized prior to estimation. This transformation eliminates scale differences and improves numerical stability, allowing us to interpret the odds ratios in terms of a one standard deviation increase in each variable.} In Table \ref{tbl:11}, we present the estimated coefficients and corresponding odds ratios from the multinomial logistic model.\footnote{For instance, an odds ratio of $0.403$ for the variable Relative income per capita in Cluster 6 means that a one standard deviation increase in income is associated with a $59.7\%$ lower likelihood of a municipality being classified in Cluster 6, relative to the non-LBP group.} These results complement the diagnosis discussed in the previous section, where it was provided a general sense of how clusters differ, but it does not account for the statistical significance of each variable. The logistic analysis, by contrast, allows us to assess the independent effect of each explanatory factor, controlling for the others, on the probability of a municipality being classified into one of the LBP clusters relative to the non-LBP group.

With respect to Category 1 (Cluster 4, rural/agricultural municipalities with moderate accessibility), the results reveal a profile strongly associated with rural economic specialization and demographic decline. Agricultural employment displays a large and highly significant positive effect (OR = 148.7), while relative income, population growth, and population density are all negatively associated with the probability of belonging to this category. These results suggest that municipalities characterized by low diversification, weak demographic dynamism, and low-density settlement patterns are substantially more likely to fall into this type of left-behind trajectory. Geographic conditions further reinforce this pattern. Distance to urban and semi-urban centres increases the likelihood of belonging to Cluster 4 (OR = 1.91), whereas market potential exerts a strong negative effect (OR = 0.003). This indicates that limited accessibility and weak integration into surrounding economic systems intensify pre-existing structural disadvantages rather than acting as isolated determinants. Overall, Cluster 4 captures a form of persistent rural disadvantage where demographic decline, agricultural specialization, and relative spatial isolation interact cumulatively.

For Category 2 (Cluster 5, structurally vulnerable municipalities with compounded disadvantages), the model points to a more complex configuration of vulnerability. The elderly dependency ratio shows a particularly strong positive association (OR = 10.8), while low income, weak population growth, and low industrial employment are also significant predictors of cluster membership. One notable feature of this category is the very large positive coefficient associated with bank branches per capita. Rather than indicating financial strength, this result likely reflects the fact that these municipalities maintain a residual service function despite broader structural fragility, possibly due to low population denominators or their role as local service centres within sparsely populated areas. Geographic variables are also highly relevant. Distance to urban centres substantially increases the likelihood of belonging to this category (OR = 189.5), whereas market potential strongly reduces it. This suggests that Cluster 5 represents territories where demographic ageing, weak economic diversification, and relative remoteness combine to produce intermediate but multidimensional forms of left-behindness.

Finally, Category 3 (Cluster 6, extreme isolation, strong decline, and absence of services) represents the most demographically fragile type of left-behind territory. The elderly dependency ratio exhibits an exceptionally strong positive effect (OR = 363.3), while population growth has the strongest negative coefficient among all categories. Together, these results indicate that municipalities experiencing intense ageing and demographic contraction are far more likely to belong to this cluster. Unlike Cluster 4, agricultural specialization and banking services are not significant defining features of this category, suggesting that its vulnerability is less associated with productive structure and more closely linked to demographic erosion and territorial marginalization. Geographic isolation nevertheless remains highly relevant, as distance to urban centres strongly increases the likelihood of belonging to this cluster (OR = 194.1), while market potential and income exert negative effects. Overall, Cluster 6 can therefore be interpreted as capturing highly marginal territories characterized by cumulative demographic decline, severe remoteness, and extremely weak territorial integration.







\subsection{Discussion}

From subsection 6.2, second-nature geographic variables show the strongest and most consistent effects across the three models. Distance to urban and semiurban centres significantly increases the probability of belonging to any type of left-behind place, with particularly high odds ratios for Clusters 5 and 6, confirming the role of spatial isolation as a key differentiating factor. Conversely, market potential (income/distance) displays a negative and significant association with all groups, indicating that municipalities with lower accessibility to markets (weaker economic interaction) are more likely to fall into the left-behind category.

These findings highlight that second-nature geography is central to understanding territorial disparities and the persistence of regional development traps in Spain.

In addition to these second-nature geographic variables, two other variables clearly differentiate the three types of left-behind places. The share of agricultural employment is strongly and positively associated with Clusters 4 and 5, revealing the persistence of a rural and low-diversified economic structure, while it loses relevance in Cluster 6. In contrast, the industrial employment share shows a consistently negative effect across all groups, suggesting that a stronger industrial base reduces the likelihood of being classified as left-behind. Together, these patterns indicate that the combination of economic structure and spatial isolation is key to understanding the heterogeneity of development traps across Spanish municipalities.

To facilitate the interpretation of these results, Figure~\ref{fig4:odds} provides a visual representation of the estimated odds ratios from the multinomial logit model. The figure displays, for each type of left-behind place, the relative weight and significance of the main socio-demographic, economic and geographical variables. This graphical summary allows for a clearer identification of the main patterns that differentiate the three clusters, confirming the robustness of the relationships discussed above.

Overall, the results underline that accessibility and market potential are decisive factors explaining why some municipalities become left-behind places. While mean comparisons offer a general overview, the multinomial logit analysis reveals distinct territorial diagnoses across clusters, emphasizing the need for targeted, cluster-specific policies that address the unique combination of demographic, economic, and accessibility challenges faced by each cluster. These findings call for differentiated policy responses tailored to the specific configurations of disadvantage. By clearly identifying the structural drivers that define Clusters 4, 5, and 6, the model provides a robust empirical foundation for designing effective and spatially balanced development strategies. Across all clusters, improving accessibility and functional connectivity must remain central pillars in any long-term approach aimed at reversing territorial marginality.

The multinomial logit results and the frontier-distance analysis jointly reveal that the three LBP clusters correspond to structurally differentiated forms of territorial disadvantage rather than to a single continuum of decline. Relative population growth emerges as the most consistent cross-cluster marker of left-behind status, displaying large negative and highly significant coefficients in the multinomial model and appearing among the variables closest to the decision boundary in all three clusters. Distance to urban or semi-urban centres also plays a common role, with positive and significant effects across clusters---reflected in odds ratios well above unity---and relatively low frontier distances, especially in Clusters 5 and 6.

Beyond these shared traits, each cluster exhibits a distinctive structural profile. Cluster 4 is primarily associated with income, market potential, population growth and agricultural specialization, while being closest to the threshold in the accessibility dimension. Cluster 5 is defined by a combination of financial accessibility, remoteness, demographic stagnation and weak market integration, confirming its intermediate and transitional nature. Cluster 6 is the most clearly demographic profile, shaped by ageing, very weak population dynamics and territorial remoteness, with the demographic dimension lying closest to the classification boundary and the socioeconomic dimension farthest from it.

Taken together, these results support the interpretation of left-behind places as heterogeneous territorial configurations driven by different combinations of demographic, socioeconomic and accessibility-related constraints.

This consistency is not coincidental. Variables that significantly affect the probability of belonging to a given cluster in the multinomial logit model are precisely those that shape the position of municipalities relative to the classification boundary. Therefore, variables with stronger effects in the logit are also expected to exhibit lower distances to the frontier, as they define the dimensions along which marginal changes can alter cluster membership.

There are some municipalities that exhibit economic and demographic factors which, without considering second-nature geography, would not be classified as LBPs; that is, their spatial conditions amplify their shortcomings.

As a result, the explicit integration of geographic variables is therefore essential for developing place-sensitive policy frameworks in Spain. Improving accessibility and functional connectivity must remain fundamental pillars of any long-term approach aimed at reversing territorial marginalization and ensuring that cohesion policy effectively addresses the challenges of all lagging places. Our recommendation aligns with the guidelines outlined in Andrés Rodríguez-Pose (2025, p. 1), which advocate for “a more dynamic and systemic approach emphasising institutional capacity, territorial sensitivity, global connections, and performance-based delivery—areas where past reforms have underdelivered.” In this way, we propose implementing so-called systemic place-based policies. Such policies do not merely respond to existing spatial structures but actively seek to reshape
them by altering the second-nature geographic preconditions for economic activity. Such policies do not merely respond to existing spatial structures, but actively seek to reshape them by altering the second-nature geographic preconditions for economic activity. They may be understood as forward-looking, place-based capacity-building strategies that contextualize local structural conditions while deliberately investing in enabling infrastructures and connectivity improvements which may exceed current demographic or economic demand, but are nevertheless necessary to unlock future development trajectories in lagging territories.








\section{Conclusions}
\label{sec:conclusions}

This paper has investigated the heterogeneous nature of left-behind places (LBPs) in Spain by developing a multidimensional framework that combines demographic, socioeconomic, and second-nature geographic variables at the municipal level. By explicitly incorporating second-nature geography variables into the clustering procedure, the analysis moves beyond approaches that treat geography as merely contextual or ex post explanatory. The analysis aimed not only to identify distinct types of territorial disadvantage, but also to explain municipalities’ proximity to vulnerability thresholds and their probability of belonging to different LBP categories. The findings show that second-nature geography constitutes a structural dimension of territorial disadvantage and plays a decisive role in shaping both the emergence and amplification of existing demographic and socioeconomic characteristics associated with vulnerability.

Our study confirmed the structural role of geographical variables by showing that their
inclusion significantly altered the classification of nearly 20\% of spanish municipalities, revealing
a "hidden" marginality obscured in purely demographic and socioeconomic analyses. The superiority of the clustering specification incorporating geographic variables is further reinforced by quantitative evidence derived from several indicators. The results of this investigation show that, while LBP clusters (CL4–CL6) account for over 71\% of Spanish municipalities, they represent only about 31\% of the total population. By contrast, NLBP clusters (CL1–CL3) concentrate nearly 69\% of the population despite comprising less than 29\% of municipalities. The analysis uncovers severe regional patterns of vulnerability, suggesting that previous assessments may have underestimated the real extent of municipal disadvantage by neglecting second-nature geographic factors. Moreover, in regions where most municipalities and population belong to LBP clusters, the results indicate  a particularly severe warning signal, that should be carefully considered in the design of economic policy. From another perspective, the results illustrate the asymmetric relationship between urbanisation and clustering: a small number of urban municipalities (largely concentrated in Cluster 1) dominate the demographic structure, whereas rural municipalities are far more numerous but fragmented across several clusters with much smaller population shares.

The extent of the emerging regional vulnerability suggests that previous estimates of vulnerability have been underestimated in previous assessments. Given that the emerging territorial vulnerability in Spain appears to have been underestimated as a result of neglecting factors associated with second-nature geography, the resulting diagnosis is particularly concerning. Regional disparities are not only likely to persist but, owing to their dependence on these structural territorial dynamics, may deepen further over time.

Methodologically, the paper also contributes by applying post-clustering sensitivity analysis based on threshold calculations to the study of municipal left-behindness, providing additional insights into the dimensions shaping municipalities’ relative proximity to vulnerability frontiers.
The results demonstrate that left-behindness in Spain is not a homogeneous phenomenon, but rather a set of differentiated territorial configurations shaped by distinct combinations of demographic decline, weak economic structures, and spatial isolation. Accessibility to urban and semiurban centres and market potential emerge as particularly important factors, both in structuring the clusters and in explaining municipalities’ position relative to vulnerability frontiers. The findings show that some municipalities only become identifiable as left-behind once second-nature geographic disadvantages are taken into account, suggesting that spatial isolation actively amplifies pre-existing structural weaknesses. Further, the combined evidence from the multinomial logit and frontier-distance analyses reinforces the robustness of the proposed typology. Therefore, it is confirmed that the dimensions exerting the strongest effects on cluster membership are also those defining the closest vulnerability thresholds. As a general conclusion, although second-nature geographic variables play an important role in marginal adjustments, being the second most important factor in determining outcomes, the multinomial analysis indicates that they remain fundamental for identifying left-behind municipalities and revealing latent forms of territorial vulnerability.

A key finding of the paper is that geography does not simply accompany territorial decline but actively amplifies pre-existing structural disadvantages. In several cases, municipalities that would not be classified as left-behind solely on the basis of demographic or socioeconomic indicators become vulnerable once second-nature geographic variables are incorporated into the analysis. This highlights the importance of considering municipalities’ functional position within broader territorial systems when analysing regional disparities and development traps.

From a policy perspective, the findings support a shift towards what we define as systemic place-based policies. The evidence suggests that cohesion policies should not only address socioeconomic disparities but also seek to transform the spatial conditions that reproduce territorial disadvantage.
Improving accessibility, functional connectivity, and integration into wider regional systems—key dimensions of second-nature geography—should remain central pillars of long-term strategies aimed at reversing territorial marginalisation. However, while addressing these second-nature geographic constraints constitutes a necessary condition for territorial recovery, it is unlikely to be sufficient in the absence of broader interventions targeting local socioeconomic structures.\footnote{At the same time, policies focused exclusively on improving connectivity are unlikely to succeed unless accompanied by measures aimed at strengthening local labour markets, productive diversification, and human capital formation.}

Systemic place-based policies, rather than simply adapting to existing territorial structures, should seek to reshape the spatial conditions constraining local development by investing in enabling infrastructures and territorial capacities capable of supporting future development trajectories, even where current demographic or economic conditions may not yet justify such investments. As a general policy recommendation, place-based policies should move beyond static compensatory approaches and adopt a more forward-looking perspective aimed at reshaping second-nature geographic conditions in lagging municipalities. Overall, the paper underlines that left-behindness is fundamentally a spatially embedded and structurally heterogeneous process, requiring systemic place-based policies capable of addressing the diverse development traps affecting municipalities in Spain.










\color{black}









%Several predictors are associated with a statistically significant decrease in the odds of being classified as a left-behind place (LBP). Specifically, increases (by one standard deviation) in relative population growth (OR = 0.077), average income relative to the national mean (OR = 0.559), net migration rate (OR = 0.605), industrial employment share (OR = 0.086), youth out-migration (OR = 0.088), and population density (OR = 0.177) reduce the probability of being classified as an LBP.

%Conversely, the probability of being classified as an LBP increases with higher values of variables such as bank presence (OR = 5.55), agricultural employment share (OR = 39.13), and the elderly dependency ratio (OR = 9.15), all of which are positively and significantly associated with LBP status.\footnote{Odds Ratios (ORs) are obtained by exponentiating standardized logistic regression coefficients. All predictors have been standardized (mean = 0, standard deviation = 1); thus, ORs represent the change in the odds of being classified as an LBP associated with a one standard deviation increase in the predictor. An OR below 1 indicates a negative association, while an OR above 1 indicates a positive association. For example, an OR of 0.60 implies a 40\% reduction in the odds; an OR of 5.0 implies that the odds are five times higher. Because variables are standardized, the magnitude of the OR reflects the impact of a one standard deviation change, allowing for comparisons of relative importance across predictors.}



%Random Forest and Rand Index

%To further evaluate and interpret the cluster structures obtained through k-means, a Random Forest classification model was applied to both typologies: one based exclusively on demographic and socio-economic variables (Model A), and another that also includes geographical connectivity indicators (Model B). The objective was twofold: to assess the internal coherence of each classification and to identify the most influential variables in explaining the differences between clusters.

%The results confirm the robustness of both models. The out-of-bag (OOB) error rate was 4.1\% for Model A and 5.9\% for Model B, suggesting a high predictive accuracy in both cases. However, the Random Forest also provides valuable insight into the role played by each variable. In Model A (without geography), the most important variables were share of agricultural employment, share of industrial employment, and external migration rate among youth, reflecting the weight of sectoral and demographic dynamics in differentiating territories.

%In contrast, in Model B (with geographical dimension), while economic and demographic variables remain important—such as share of industrial employment and share of agricultural employment—two geographical connectivity indicators stand out: distance to the nearest urban centre (dist\_1\_std, MeanDecreaseGini = 340.5) and market potential accessibility (PM\_renda\_std, MeanDecreaseGini = 420.2). The appearance of these variables among the top-ranking features supports the argument that spatial connectivity plays a significant role in shaping territorial disadvantage. In Figure \ref{RF} the comparison of the role played by each variable in each model is represented.

%Thus, the Random Forest analysis not only reinforces the internal consistency of the clusters but also highlights how the inclusion of territorial connectivity variables changes the explanatory structure of the typology. This finding underscores the importance of geography, not just as a background condition but as a central explanatory dimension for understanding local development trajectories.

%MATERIAL INTERESANTE
%Finally, to complement this analysis, a Rand index was calculated to measure the similarity between the two cluster classifications. The result (Rand Index = 0.56) suggests a moderate correspondence, indicating that the inclusion of geographical factors leads to a substantively different territorial segmentation. This reinforces the analytical relevance of incorporating spatial variables when identifying and characterizing left behind places.

%There are significant overlaps, but also notable differences.
%Many municipalities likely fell into similar clusters, but some were reassigned when the additional variables were included.
%This suggests that the two additional variables did influence the clustering results, although not in a radical way.
%Importantly, the most substantial changes in cluster membership affect municipalities that could be considered left behind places (LBP), indicating that the inclusion of these variables helps to better differentiate vulnerable or lagging areas.

%Clustering
%Comparación de clustering
%Logit
%Although the difference in silhouette scores (0.34 vs. 0.31) is relatively small and unlikely to be statistically significant, it does reflect meaningful differences in other dimensions of the clustering. In particular, the solution with the lower silhouette value (0.31), which incorporates two additional explanatory variables, yields more coherent and interpretable cluster profiles. This suggests that the inclusion of these variables enhances the substantive relevance of the clustering, despite the marginal decrease in overall compactness.

\pagebreak
\clearpage

% \theendnotes

\pagebreak
\clearpage


\bibliography{Bib.files/territorial_v5}

\pagebreak
\clearpage

\appendix
\section{Appendix A: Distance metrics}

\label{sec:appendix_A}

After testing the three metrics, we ultimately chose to use \textbf{Euclidean distance} for the clustering process. The \textbf{Canberra distance}, although valuable in cases involving heterogeneous units of measurement like ours, did not yield satisfactory results due to the numerical and relatively homogeneous nature of the variables. On the other hand, the \textbf{Mahalanobis distance} produced coherent results and took into account the correlation structure among variables. However, it did not align well with the conceptual aim of our analysis.

Our objective was to group municipalities based on their absolute similarity across a set of socioeconomic indicators. In other words, to cluster them based on proximity in the multidimensional space defined by the variables. The Euclidean distance directly reflects this idea by giving equal weight to all variables and focusing on overall closeness in magnitude. In contrast, Mahalanobis distance tends to form clusters by adjusting for the statistical distribution of the data. For instance, two municipalities with very different values in highly variable indicators could still be considered close under this metric. This means it emphasizes statistical patterns over raw similarity, which can lead to groupings that do not correspond to real-world similarities. Therefore, despite the higher silhouette values achieved with Mahalanobis distance in some configurations, we considered Euclidean distance to be a more conceptually faithful choice for our study’s purposes.

Nonetheless, it is important to acknowledge that Euclidean distance has two limitations, it is particularly sensitive to outliers and to differences in scale among variables. To mitigate these effects and ensure the robustness of the clustering process, a treatment of atypical data and standardization have been carried out.

First, a winsorization procedure was implemented to reduce the influence of extreme values without removing any observations. Specifically, the Stata command \texttt{winsor2} was applied, which caps the lower and upper tails of each variable distribution at the 5th and 95th percentiles, respectively. This approach helps limit the impact of outliers while preserving the integrity of the dataset. An exception was made for the variable \texttt{dist}, for which the lower tail (5th percentile) was not modified. This decision was based on the fact that values close to zero are conceptually meaningful in this context, as they correspond to municipalities located in urban or semi-urban areas. Altering these values would have distorted the information carried by the variable, hence they were retained in their original form.

Second, all variables were standardized. As previously discussed, the algorithm we will use based on the Euclidean distance, \(k\)-means, can assign disproportionate weight to variables with larger variances or magnitudes. To avoid this bias and ensure that all variables contribute equally to the distance computation, the data were transformed to have zero mean and unit standard deviation.

\begin{comment}
\begin{itemize}
    \item \textbf{Euclidean distance} (Equation~\ref{eq:eq23}) is the most intuitive metric and measures the straight-line distance between two points in a multidimensional space. It is especially appropriate when we are interested in clustering based on the absolute similarity of variables.

    \item \textbf{Mahalanobis distance} (Equation~\ref{eq:eq24}) accounts for correlations among variables and scales the distances accordingly. This is useful when variables have different variances or when multicollinearity is present.

    \item \textbf{Canberra distance} (Equation~\ref{eq:eq25}) is less sensitive to absolute differences and gives more weight to small changes when values are near zero. It is often used when variables differ significantly in scale or units, even after normalization.
\end{itemize}
\end{comment}


\begin{comment}
    Let \(m_i\) and \(m_j\) be any two municipalities, \(p\) the number of variables, \(x_{i}\) and \(x_{j}\) the corresponding vectors of values, \(x_{il}\) and \(x_{jl}\) (with \(1 \leq l \leq p\)) the individual values for each variable, and \(\mathbf{S}\) the variance-covariance matrix of the dataset. The metrics are defined as follows:
\begin{equation}
    \label{eq:eq23}
    d_{E}(m_i, m_j) = \sqrt{\sum^p_{l=1}(x_{il} - x_{jl})^2}
\end{equation}
\begin{equation}
    \label{eq:eq24}
    d_{M}(m_i, m_j) = \sqrt{(x_i - x_j)^{T} \, \mathbf{S}^{-1} \, (x_i - x_j)}
\end{equation}
\begin{equation}
    \label{eq:eq25}
    d_{C}(m_i, m_j) = \sum^p_{l=1}\frac{|x_{il} - x_{jl}|}{|x_{il} + x_{jl}|}
\end{equation}



Partitioning methods start from a simple premise: the dataset, composed of \(n\) municipalities, is to be divided into \(K\) clusters, with \(K < n\). The number of clusters \(K\) must be specified in advance, and in this study it has been determined by combining two criteria, the elbow method and the maximization of the silhouette coefficient.

Each municipality \(m_i\) is assigned to a cluster \(k \in \{1, 2, \dots, K\}\) through a classification function \(C\), so that \(C(m_i) = k\). The goal is to find the optimal partitioning \(C\) that minimizes the Population intra-cluster dissimilarity, defined as:
\begin{equation}
    \label{eq:eq26}
    W(C) = \sum_{k=1}^{K} \sum_{C(i)=k} \sum_{C(j)=k} d(x_i, x_j)
\end{equation}

Depending on the metric used, different clustering procedures can be applied. In the case of the Euclidean distance, the standard method is \textit{\(k\)-means} clustering. Under this method, equation~\ref{eq:eq26} becomes:
\[
W(C) = \sum_{k=1}^{K} \sum_{C(i)=k} \sum_{l=1}^p (x_{il} - x_{jl})^2
\]

By applying a known identity relating pairwise squared distances and means, this can be rewritten in a simplified and computationally convenient form:
\[
W(C) = \sum_{k=1}^{K} \sum_{C(i)=k} \sum_{l=1}^p (x_{il} - \Bar{x}_k^l)^2
\]
where \(\Bar{x}_k = (\Bar{x}_k^1, \dots, \Bar{x}_k^p)\) is the vector of means for cluster \(k\). This formulation shows that the optimization problem consists of minimizing the Population squared distance of each observation from the centroid of its assigned cluster:
\[
\min_{C, \{m_{kl}\}} W(C, \{m_{kl}\}) = \sum_{k=1}^{K} \sum_{C(i)=k} \sum_{l=1}^p (x_{il} - m_{kl})^2
\]

This is a combinatorial optimization problem that is computationally infeasible to solve exactly when \(n\) is large. Therefore, iterative approximation algorithms are commonly used. The most widely known is Lloyd's algorithm, which alternates between assigning points to the nearest cluster mean and updating the means based on the new assignments. This process continues until convergence is reached.

It is important to mention here that the \(k\)-means algorithm is, in fact, a specific case of a more general method known as \textit{\(k\)-medoids}. While \(k\)-means uses the arithmetic mean as the cluster representative, \(k\)-medoids selects an actual observation (called a medoid) to represent each cluster. This makes the method more robust and better suited for arbitrary distance metrics, such as Mahalanobis or Canberra, for which computing a meaningful mean is not straightforward.

However, despite its flexibility, we did not adopt the \(k\)-medoids method in this study. Since we are working exclusively with the Euclidean distance, and because \(k\)-means is computationally more efficient and widely accepted in the economic literature for large-scale regional classifications, it was deemed more appropriate for our analysis.

Once the clustering procedures have been applied, it is essential to evaluate the quality of the resulting classifications. Among the available validation tools, we use the \textit{silhouette coefficient}, a widely accepted internal measure that assesses how well each observation has been classified within its assigned cluster.
\end{comment}


\section{Appendix B: Determination of the number of clusters}
\label{sec:appendix_B}

\subsection{Silhouette index}




The silhouette coefficient provides a measure of how well each observation is assigned to its cluster relative to other clusters. For an individual observation \(i\), the silhouette coefficient is defined as:

\[
s(i) = \frac{b(i) - a(i)}{\max\{a(i), b(i)\}}
\]

Here, \(a(i)\) denotes the average dissimilarity between observation \(i\) and the other members of its own cluster:
\[
a(i) = \frac{1}{|C(i)| - 1} \sum_{j \in C(i), j \ne i} d(x_i, x_j)
\]

Meanwhile, \(b(i)\) represents the minimum average dissimilarity between observation \(i\) and all points in the nearest external cluster, i.e., the one with the lowest average distance to \(i\):
\[
b(i) = \min_{k \ne C(i)} \left( \frac{1}{|C_k|} \sum_{j \in C_k} d(x_i, x_j) \right)
\]

The silhouette coefficient \(s(i)\) ranges from \(-1\) to \(1\). A value close to 1 indicates that the observation is well classified, being significantly closer to members of its own cluster than to those of any other. A value near 0 suggests that the observation lies at the boundary between two clusters, making its assignment ambiguous. Finally, values approaching \(-1\) imply a likely misclassification, as the observation is closer to a different cluster than its own.

To obtain a global assessment of clustering performance, we compute the average silhouette value across all observations:
\begin{equation}
    \label{eq:eq28}
    \bar{s} = \frac{1}{n} \sum_{i=1}^{n} s(i)
\end{equation}

This average value \(\bar{s}\) serves as an overall indicator of the cohesion and separation of the clusters. In general, a mean silhouette above 0.5 suggests well-separated and well-defined clusters. Values between 0.2 and 0.5 are considered acceptable, providing meaningful structure, while those below 0.2 indicate a weak or ambiguous clustering.

To determine the appropriate number of clusters, the silhouette index was computed for alternative values of $K$. Figure \ref{fig:silhouette}  presents the average silhouette values obtained for the two clustering specifications considered in the analysis. The first specification includes only socio-economic variables, while the second additionally incorporates geographical indicators capturing spatial accessibility and market potential.

In both cases, the silhouette index reaches its maximum when $K=6$, indicating that this number of clusters provides the most appropriate balance between within-cluster cohesion and between-cluster separation. The consistency of this result across both specifications supports the choice of a six-cluster solution for the subsequent analysis.

\begin{figure}[ht]
\centering
\caption{Average silhouette values for alternative numbers of clusters. In both specifications the maximum value is obtained for $K=6$.}
\label{fig:silhouette}
\begin{minipage}{0.48\textwidth}
\centering
\includegraphics[width=\textwidth]{Figures/siluete_sin geog.png}
\caption*{(a) Without geographical and functional connectivity variables}
\end{minipage}
\hfill
\begin{minipage}{0.48\textwidth}
\centering
\includegraphics[width=\textwidth]{Figures/siluete_geo.png}
\caption*{(b) Including geographical and functional connectivity variables}
\end{minipage}

\end{figure}

\subsection{Elbow criterion}

The elbow method is a commonly used heuristic to determine the appropriate number of clusters in partitioning algorithms. The method examines the evolution of WCSS, shown in Equation \ref{eq:within_cluster_variance}, as the number of clusters $K$ increases. Since WCSS decreases monotonically with $K$, the optimal number of clusters is identified at the point where the reduction in WCSS begins to slow down, forming an ``elbow'' in the curve. This point indicates that adding further clusters provides only marginal improvements in within-cluster homogeneity.

Figure \ref{fig:elbow}  present the elbow plots for the two clustering specifications considered in the analysis. In both cases,
the plot also supports the choice of six clusters, as the reduction in within-cluster sum of squares becomes noticeably smaller beyond $K=6$, indicating diminishing returns from increasing the number of clusters.

\begin{figure}[ht]
\centering
\caption{Elbow plots for alternative numbers of clusters. In both specifications the reduction in WCSS becomes noticeably smaller beyond $K=6$, supporting the selection of a six-cluster solution.}
\label{fig:elbow}

\begin{minipage}{0.48\textwidth}
\centering
\includegraphics[width=\textwidth]{Figures/colze_singeo.png}
\caption*{(a) Without geographical and functional connectivity variables}
\end{minipage}
\hfill
\begin{minipage}{0.48\textwidth}
\centering
\includegraphics[width=\textwidth]{Figures/colze_congeo.png}
\caption*{(b) Including geographical and functional connectivity variables}
\end{minipage}

\end{figure}


%%%%%%%%%%%%%%%%%%%%%%%%


\pagebreak
\clearpage

\section{Appendix C: The k-Means Algorithm}
\label{sec:appendix_c}
As explained in \citep{james2013introduction}, the goal is to find the optimal partitioning $C$ that minimizes the total within-cluster variance, defined as the sum of squared Euclidean distances from each observation ($x_i$) to the centroid ($\mu_k$) of its assigned cluster $k$. The function known as within-cluster sum of squares (WCSS), namely:
\begin{equation}
 \label{eq:within_cluster_variance}
 W(C) = \sum_{k=1}^{K} \sum_{C(i)=k} \|x_i - \mu_k\|^2
\end{equation}

{\color{purple}

Expanding the squared Euclidean distance into its component-wise representation, the objective function can be written as:
\[
W(C) = \sum_{k=1}^{K} \sum_{C(i)=k} \sum_{l=1}^p (x_{il} - \Bar{x}_k^l)^2
\]
where \(\Bar{x}_k = (\Bar{x}_k^1, \dots, \Bar{x}_k^p)\) is the vector of means for cluster \(k\). This formulation shows that the optimization problem consists of minimizing the population squared distance of each observation from the centroid of its assigned cluster:
\[
\min_{C, \{m_{kl}\}} W(C, \{m_{kl}\}) = \sum_{k=1}^{K} \sum_{C(i)=k} \sum_{l=1}^p (x_{il} - m_{kl})^2
\]
This is a combinatorial optimization problem that becomes computationally infeasible to solve exactly when \(n\) is large. Therefore, iterative approximation algorithms—such as Lloyd’s algorithm—are used to obtain a local optimum. Moreover, since the algorithm may converge to local minima, it is standard practice to run the procedure with multiple random initializations and retain the best solution.
}


It is important to mention that the $k$-means algorithm is, in fact, a specific case of a more general method known as $k$-medoids. While $k$-means uses the arithmetic mean as the cluster representative, $k$-medoids selects an actual observation (called a medoid) to represent each cluster. This makes the method more robust to outliers. However, despite its flexibility, we did not adopt the $k$-medoids method in this study. Since we are working exclusively with the Euclidean distance, and because $k$-means is computationally more efficient and widely accepted in the economic literature for large-scale regional classifications, it was deemed more appropriate for our analysis.  {\color{red} (EMILI: quiz\'as justificar m\'as los m\'etodos indicando que son datos de municipios, con mucha variabilidad, etc.)}

\begin{sidewaystable}
\centering
\begin{threeparttable}
\scriptsize
\caption{Description of variables}
\label{tbl:1}
\begin{tabular}{ll>{\raggedright}p{4.3cm}>{\raggedright\arraybackslash}p{5.2cm}>{\centering\arraybackslash}p{2cm}}
\toprule
Type & Variable & Definition & Description & Source \\
\midrule

\textbf{Socio-demographic} \\
& \(\text{DEPEND}_{i}\) & Elderly dependency ratio & Ratio of population aged 65+ to population aged 16–64 (2022) & INE \\
& \(\text{ShFPOP16}_{i}\) & Share of foreign-born under 16 & Foreign population under 16 divided by total foreign population (2022) & INE \\
& \(\text{NETMIGR}^{i}\) & Net migration rate & Net migration flow (immigrants – emigrants) divided by total population (2022) & INE \\
& \(\Delta\text{POP}_{i}\) & Relative population growth & Difference in growth rate between the municipality and its autonomous community (2000–2022) & INE \\
& \(\text{DENSPOP}_{i}\) & Population density & Population divided by surface area (inhabitants/km\(^2\)) in 2022 & INE \\

\midrule
\textbf{Economic and financial} \\
& \(\text{INCOME}_{PCi}\) & Relative income per capita & Municipal income per capita divided by national income per capita (2022) & INE \\
& \(\text{ShINDUST}_{i}\) & Industry employment share & Share of labour contracts in the industrial sector (2022) & SEPE \\
& \(\text{ShAGRICULT}_{i}\) & Agriculture employment share & Share of labour contracts in the agricultural sector (2022) & SEPE \\
& \(\text{ShBRANCHES}_{i}\) & Bank branches per capita & Number of bank branches per inhabitant (2022) & Bank of Spain \\

\midrule
\textbf{Geographical and distance-related} \\
& \(\text{DISTANCE}_{i}\) & Composite distance index & Harmonic mean of distances to the nearest urban and semi-urban centres & Own \\
& \(PM_{it}\) & Market potential & Sum of income in surrounding municipalities weighted by the inverse of squared distance & Own \\

\bottomrule
\end{tabular}

\begin{tablenotes}
\footnotesize
\item[] \textit{Source:} INE (Instituto Nacional de Estadística), SEPE (Servicio Público de Empleo Estatal), Bank of Spain. Distance and market potential variables are based on authors’ calculations using geospatial data.
\end{tablenotes}

\end{threeparttable}
\end{sidewaystable}



\clearpage
\pagebreak





\clearpage
\pagebreak


\begin{table}[htbp]
\centering
\caption{MANOVA results without geographical variables}
\label{tab:manova_nogeo}
\footnotesize
\begin{tabular}{lcccc}
\toprule
Test & Value & Approx. F & Num df & Den df \\
\midrule
Wilks' Lambda & 0.0027 & 2219.3 & 45 & 36160 \\
Pillai's Trace & 3.3700 & 1857.8 & 45 & 40435 \\
\bottomrule
\end{tabular}
\end{table}




\begin{table}[htbp]
\centering
\caption{MANOVA results with geographical variables}
\label{tab:manova_geo}
\footnotesize
\begin{tabular}{lcccc}
\toprule
Test & Value & Approx. F & Num df & Den df \\
\midrule
Wilks' Lambda & 0.0033 & 1656.2 & 55 & 37409 \\
Pillai's Trace & 3.3037 & 1431.5 & 55 & 40425 \\
\bottomrule
\end{tabular}
\end{table}

\clearpage
\pagebreak

\begin{table}[!ht]
    \centering
    \caption{Cluster Means (Excluding geographical variables)}
    \label{tbl:2}
    \begin{adjustbox}{max width=\textwidth}
    \begin{tabular}{lrrrrrr}
        \toprule
        \textbf{Variable} & \textbf{Cluster 1} & \textbf{Cluster 2} & \textbf{Cluster 3} & \textbf{Cluster 4} & \textbf{Cluster 5} & \textbf{Cluster 6} \\
        \midrule
        \multicolumn{7}{l}{\textit{Socio-demographic}} \\
        Elderly dependency ratio                    & 0.29 & 0.40 & 0.36 & 0.44 & 0.57 & 0.74 \\
        Share of foreign-born under 16              & 0.15 & 0.14 & 0.17 & 0.13 & 0.15 & 0.03 \\
        Net migration rate                          & 0.01 & 0.01 & 0.02 & 0.01 & 0.01 & 0.01 \\
        Relative population growth (2000–2022)      & 0.16 & 0.01 & 0.23 & -0.20 & -0.32 & -0.34 \\
        Population density (inhabitants/km\(^2\))   & 2,530.99  & 119.92 & 160.11 & 44.74 & 12.66 & 7.99 \\
        \addlinespace

        \multicolumn{7}{l}{\textit{Economic and Financial}} \\
        Relative income per capita                  & 1.04  & 1.02 & 1.00 & 0.81 & 0.95 & 1.01 \\
        Industry employment share                   & 0.12  & 0.50 & 0.07 & 0.05 & 0.08 & 0.02 \\
        Agriculture employment share                & 0.03  & 0.07 & 0.06 & 0.54 & 0.21 & 0.04 \\
        Bank branches per capita                    & 0.0003  & 0.0005 & 0.0003 & 0.0005 & 0.0031 & 0.0001 \\
        \midrule

        \textbf{Number of municipalities}           & 222  & 785 & 2,260 & 1,722 & 523 & 2,585 \\
        \bottomrule
    \end{tabular}
    \end{adjustbox}
    \vspace{0.4cm}
    \raggedright
    \textit{Source: Authors’ own elaboration.}
\end{table}



\begin{table}[htbp]
\centering
\caption{Cluster Means (Including geographical variables)}
\label{tab:cluster_means}

\resizebox{\textwidth}{!}{

\begin{tabular}{lcccccc}
\toprule
Variable & Cluster 1 & Cluster 2 & Cluster 3 & Cluster 4 & Cluster 5 & Cluster 6 \\
\midrule

\multicolumn{7}{l}{\textit{Socio-demographic}} \\
\midrule
Elderly dependency ratio & 0.293 & 0.395 & 0.309 & 0.428 & 0.561 & 0.749 \\
Share under 16 & 0.154 & 0.144 & 0.135 & 0.144 & 0.147 & 0.045 \\
Net migration rate & 0.004 & 0.009 & 0.018 & 0.008 & 0.009 & 0.007 \\
Relative population growth & 0.115 & -0.009 & 0.461 & -0.154 & -0.309 & -0.343 \\Population density & 2563.87 & 118.42 & 206.67 & 66.98 & 12.32 & 6.55 \\

\midrule
\multicolumn{7}{l}{\textit{Economic and Financial}} \\
\midrule
Relative income per capita & 1.041 & 1.023 & 1.080 & 0.823 & 0.943 & 1.025 \\
Industry employment share & 0.130 & 0.503 & 0.079 & 0.063 & 0.082 & 0.017 \\
Agriculture employment share & 0.030 & 0.062 & 0.068 & 0.358 & 0.218 & 0.058 \\
Bank branches per capita & 0.00030 & 0.00043 & 0.00021 & 0.00046 & 0.00302 & 0.00004 \\

\midrule
\multicolumn{7}{l}{\textit{Geographical and distance-related factors}} \\
\midrule
Distance to urban and semiurban centers & 0.11 & 6.36 & 4.57 & 9.04 & 17.73 & 17.51 \\
Market potential ($PM\_renda$) & 139,243,952 & 98,588,152 & 117,940,928 & 46,838,040 & 62,413,868 & 88,904,232 \\


\bottomrule
\end{tabular}

}

\vspace{0.2cm}
\footnotesize{Source: Authors' own elaboration.}

\end{table}

\clearpage
\pagebreak













\clearpage
\pagebreak
\begin{table}[ht]
\centering
\caption{Cluster Typology of Spanish Municipalities and Left-Behind Status}
\label{tbl:5}

\small

\begin{adjustbox}{max width=\textwidth}
\begin{tabular}{clp{8cm}cc}
\toprule
\textbf{Cluster} & \textbf{Name} & \textbf{Main Characteristics} & \textbf{Vulnerability Status} & \textbf{\% Municipalities} \\
\midrule

\textbf{1} & \textbf{Urban core / highly connected}
& Dense urban municipalities with the highest income and market potential, very high accessibility, limited agricultural activity, and relatively young populations.
& \textbf{NO}
& 2.7\% \\

\textbf{2} & \textbf{Industrial and accessible}
& Industrial municipalities with good accessibility, relatively high income levels, moderate density, and stable demographic structures.
& \textbf{NO}
& 9.0\% \\

\textbf{3} & \textbf{Dynamic semi-central}
& Growing semi-central municipalities with diversified economic structures, moderate density, good accessibility, and favourable demographic dynamics.
& \textbf{NO}
& 17.2\% \\

\textbf{4} & \textbf{Rural/agricultural with moderate accessibility}
& Predominantly rural and agricultural municipalities with moderate accessibility, ageing tendencies, lower income levels, and limited service provision, showing signs of emerging vulnerability.
& \textbf{Borderline}
& 31.9\% \\

\textbf{5} & \textbf{Structurally vulnerable rural areas}
& Structurally fragile rural municipalities characterised by low density, demographic decline, high dependency ratios, weak accessibility, and limited integration into major economic centres.
& \textbf{YES}
& 6.9\% \\

\textbf{6} & \textbf{Critical rural periphery}
& Extremely remote and sparsely populated municipalities facing severe demographic decline, ageing, very limited services, weak human capital, and persistent territorial isolation.
& \textbf{YES (critical)}
& 32.4\% \\

\bottomrule
\end{tabular}
\end{adjustbox}

\smallskip
\textit{Source: Authors' own elaboration.}

\end{table}



\clearpage
\pagebreak

\begin{table}[H]
\centering
\scriptsize
\caption{Reassignment of municipalities between clusters according to the inclusion of geographical variables (transition matrix)}
\label{tab:reassign_clusters}
\setlength{\tabcolsep}{5pt}
\renewcommand{\arraystretch}{1.1}
\begin{tabular}{l *{7}{r} r}
\toprule
\textbf{Without taking geography into account}
& \multicolumn{7}{c}{\textbf{Taking geography into account}}  \\
\cmidrule(lr){2-8}
 & \textbf{CL1} & \textbf{CL2} & \textbf{CL3} & \textbf{CL4} & \textbf{CL5} & \textbf{CL6} & \textbf{Total} \\
\midrule
CL1 & 212 & 0 & 7 & 3 & 0 & 0 & 222 \\
CL2 & 2 &695 & 32 & 35 & 8 & 13 & 785 \\
CL3 & 1 & 33 & 1,216 & 751 & 26 & 233 &  2,260 \\
CL4 & 0 & 2 & 42 & 1,578 & 22 & 78 &  1,722 \\
CL5 & 0 & 3 & 3 & 15 & 500 & 2 &  523 \\
CL6 & 0 & 3 & 89 & 196 & 2 & 2,295 &  2,585 \\
\midrule
\textbf{Total}
& 215 & 736 & 1,389 & 2,578 & 558 & 2,621 &  8,097 \\
\bottomrule
\end{tabular}

\vspace{2mm}
\footnotesize\emph{Note: The category ``No change'' includes municipalities whose cluster assignment remains unchanged after incorporating geographical variables.}

\footnotesize\emph{Source: Authors' own elaboration.}
\end{table}



\clearpage
\pagebreak

\begin{table}[ht]
\centering
\caption{Spatial autocorrelation of cluster assignments (Moran's I)}
\label{tab:moran}
\begin{tabular}{lccc}
\hline
Specification & Moran's $I$ & Z-score & p-value \\
\hline
Without geographical variables & 0.334 & 49.76 & $<0.001$ \\
With geographical variables & 0.419 & 62.35 & $<0.001$ \\
\hline
\end{tabular}

\end{table}

%%%%%%%%%%%%%%%%%%%%%%%



%%%%%%%%%%%%%%%%%%%%


%%%%%%%%%%%%%%%%%%%%%%%%%%

\clearpage
\pagebreak

\sisetup{group-separator = {,}, group-minimum-digits = 3}
\newcommand{\pnum}[1]{\num[round-mode=places,round-precision=2,minimum-decimal-digits=2]{#1}}

% --- ESTIL MÉS LLEGIBLE ---
\footnotesize
\setlength{\tabcolsep}{3pt}
\renewcommand{\arraystretch}{1.05}





\sisetup{group-separator = {,}, group-minimum-digits = 3}
% \newcommand{\pnum}[1]{\num[round-mode=places,round-precision=2,minimum-decimal-digits=2]{#1}}

\newcommand{\yourcommandname}[1]{\num[round-mode=places,round-precision=2]{#1}}

% --- ESTIL COMPRIMIT PER A PÀGINA VERTICAL ---
\scriptsize                         % més menut que \footnotesize
\setlength{\tabcolsep}{2pt}         % menys espai entre columnes
\renewcommand{\arraystretch}{0.92}  % menys espai vertical a les files


\pagebreak
\clearpage



\begin{table}%[p]
\centering
\caption{Regional composition of clusters (100\% = each cluster; population and municipalities)}
\label{tbl:7}
\footnotesize
\begin{adjustbox}{max totalsize={\textwidth}{\textheight},center}
\begin{tabular}{llrrrrrr}
\toprule
\textbf{CCAA} &  & \textbf{CL1} & \textbf{CL2} & \textbf{CL3} & \textbf{CL4} & \textbf{CL5} & \textbf{CL6} \\
\midrule

\multirow{2}{*}{\textbf{Andalusia}}
 & pct\_mun (\%) & \pnum{11.62790698} & \pnum{2.197802198} & \pnum{3.019410496} & \pnum{25.30027121} & \pnum{3.220035778} & \pnum{0.571864278} \\
 & pct\_pop (\%) & \pnum{12.24696492} & \pnum{3.541525591} & \pnum{8.424530523} & \pnum{37.2350914}  & \pnum{2.563145236} & \pnum{0.70335754} \\
\midrule

\multirow{2}{*}{\textbf{Aragon}}
 & pct\_mun (\%) & \pnum{0.465116279} & \pnum{7.554945055} & \pnum{3.666427031} & \pnum{4.998062766} & \pnum{32.55813953} & \pnum{11.89477697} \\
 & pct\_pop (\%) & \pnum{0.00132758}  & \pnum{6.57861262}  & \pnum{8.243620009} & \pnum{1.26299652}  & \pnum{29.48109044} & \pnum{5.297329111} \\
\midrule

\multirow{2}{*}{\textbf{Asturias}}
 & pct\_mun (\%) & \pnum{0.930232558} & \pnum{1.236263736} & \pnum{0.431344357} & \pnum{1.084850833} & \pnum{1.431127013} & \pnum{0.953107129} \\
 & pct\_pop (\%) & \pnum{1.740968999} & \pnum{1.688193254} & \pnum{3.034506398} & \pnum{1.757108899} & \pnum{1.217030908} & \pnum{9.619092578} \\
\midrule

\multirow{2}{*}{\textbf{Balearic Islands}}
 & pct\_mun (\%) & \pnum{0.930232558} & \pnum{0} & \pnum{2.372393961} & \pnum{1.123595506} & \pnum{0.178890877} & \pnum{0.07624857} \\
 & pct\_pop (\%) & \pnum{2.364586979} & \pnum{0} & \pnum{4.111223353} & \pnum{2.121374909} & \pnum{0.163814386} & \pnum{0.080160382} \\
\midrule

\multirow{2}{*}{\textbf{Canary Islands}}
 & pct\_mun (\%) & \pnum{1.860465116} & \pnum{0} & \pnum{1.150251618} & \pnum{2.518403719} & \pnum{0} & \pnum{0.07624857} \\
 & pct\_pop (\%) & \pnum{3.195874936} & \pnum{0} & \pnum{4.315859354} & \pnum{8.038203908} & \pnum{0} & \pnum{0.943072334} \\
\midrule

\multirow{2}{*}{\textbf{Cantabria}}
 & pct\_mun (\%) & \pnum{2.325581395} & \pnum{1.785714286} & \pnum{2.156721783} & \pnum{1.123595506} & \pnum{0.178890877} & \pnum{0.914982844} \\
 & pct\_pop (\%) & \pnum{1.30812222}  & \pnum{3.522124108} & \pnum{1.741825677} & \pnum{0.303642762} & \pnum{0.384934865} & \pnum{2.089534395} \\
\midrule

\multirow{2}{*}{\textbf{Castile--La Mancha}}
 & pct\_mun (\%) & \pnum{0.465116279} & \pnum{5.357142857} & \pnum{6.25449317}  & \pnum{12.39829523} & \pnum{21.82468694} & \pnum{13.30537552} \\
 & pct\_pop (\%) & \pnum{0.178270669} & \pnum{4.97927363}  & \pnum{6.101935958} & \pnum{8.454069151} & \pnum{17.15015065} & \pnum{7.130748787} \\
\midrule

\multirow{2}{*}{\textbf{Castile and León}}
 & pct\_mun (\%) & \pnum{2.790697674} & \pnum{17.58241758} & \pnum{16.534867} & \pnum{13.32816738} & \pnum{22.71914132} & \pnum{53.83149066} \\
 & pct\_pop (\%) & \pnum{3.894141415} & \pnum{6.341340786} & \pnum{6.610556595} & \pnum{2.98604232}  & \pnum{30.29032193} & \pnum{37.79447062} \\
\midrule

\multirow{2}{*}{\textbf{Catalonia}}
 & pct\_mun (\%) & \pnum{34.41860465} & \pnum{25.96153846} & \pnum{26.24011503} & \pnum{6.509104998} & \pnum{4.830053667} & \pnum{4.689287076} \\
 & pct\_pop (\%) & \pnum{26.48039471} & \pnum{23.95901984} & \pnum{15.21374975} & \pnum{2.278446892} & \pnum{4.793162631} & \pnum{7.362953335} \\
\midrule

\multirow{2}{*}{\textbf{Valencian Community}}
 & pct\_mun (\%) & \pnum{20.46511628} & \pnum{10.98901099} & \pnum{9.777138749} & \pnum{7.438977141} & \pnum{0.894454383} & \pnum{3.202439954} \\
 & pct\_pop (\%) & \pnum{10.58513387} & \pnum{22.43475779} & \pnum{9.683492981} & \pnum{10.19561638} & \pnum{0.881298243} & \pnum{4.561784802} \\
\midrule

\multirow{2}{*}{\textbf{Extremadura}}
 & pct\_mun (\%) & \pnum{0} & \pnum{0.961538462} & \pnum{0.143781452} & \pnum{11.85586982} & \pnum{1.967799642} & \pnum{2.13495997} \\
 & pct\_pop (\%) & \pnum{0} & \pnum{0.498550173} & \pnum{0.944239021} & \pnum{6.594880648} & \pnum{3.390899908} & \pnum{3.684618709} \\
\midrule

\multirow{2}{*}{\textbf{Galicia}}
 & pct\_mun (\%) & \pnum{0.930232558} & \pnum{4.120879121} & \pnum{1.797268152} & \pnum{6.741573034} & \pnum{1.252236136} & \pnum{2.783072817} \\
 & pct\_pop (\%) & \pnum{2.721407008} & \pnum{8.649739816} & \pnum{7.155251357} & \pnum{7.757969471} & \pnum{2.054336595} & \pnum{17.83882705} \\
\midrule

\multirow{2}{*}{\textbf{Community of Madrid}}
 & pct\_mun (\%) & \pnum{7.906976744} & \pnum{0.824175824} & \pnum{9.345794393} & \pnum{0.658659434} & \pnum{0} & \pnum{0.343118567} \\
 & pct\_pop (\%) & \pnum{26.9281793}  & \pnum{1.324207842} & \pnum{13.17929665} & \pnum{0.415265592} & \pnum{0} & \pnum{0.234043792} \\
\midrule

\multirow{2}{*}{\textbf{Region of Murcia}}
 & pct\_mun (\%) & \pnum{0.465116279} & \pnum{0.274725275} & \pnum{0.287562904} & \pnum{1.472297559} & \pnum{0} & \pnum{0} \\
 & pct\_pop (\%) & \pnum{0.216010421} & \pnum{1.122304081} & \pnum{1.0623462}   & \pnum{9.881735671} & \pnum{0} & \pnum{0} \\
\midrule

\multirow{2}{*}{\textbf{Navarre}}
 & pct\_mun (\%) & \pnum{3.720930233} & \pnum{8.516483516} & \pnum{5.391804457} & \pnum{1.007361488} & \pnum{5.366726297} & \pnum{2.706824247} \\
 & pct\_pop (\%) & \pnum{1.507649368} & \pnum{4.76706519}  & \pnum{1.551423004} & \pnum{0.42118166}  & \pnum{4.102595271} & \pnum{1.604433804} \\
\midrule

\multirow{2}{*}{\textbf{Basque Country}}
 & pct\_mun (\%) & \pnum{10.23255814} & \pnum{9.89010989} & \pnum{9.705248023} & \pnum{0.309957381} & \pnum{1.610017889} & \pnum{0.152497141} \\
 & pct\_pop (\%) & \pnum{5.870801446} & \pnum{8.777321554} & \pnum{7.911649705} & \pnum{0.031437798} & \pnum{1.570997329} & \pnum{0.138862918} \\
\midrule

\multirow{2}{*}{\textbf{La Rioja}}
 & pct\_mun (\%) & \pnum{0.465116279} & \pnum{2.747252747} & \pnum{1.725377426} & \pnum{2.130956993} & \pnum{1.967799642} & \pnum{2.363705681} \\
 & pct\_pop (\%) & \pnum{0.760166158} & \pnum{1.815963721} & \pnum{0.714493465} & \pnum{0.264936025} & \pnum{1.956221618} & \pnum{0.916709838} \\
\midrule

\multicolumn{8}{l}{\textbf{Totals (Spain)}} \\
\quad \% Municipalities & & \num{2.66} & \num{8.99} & \num{17.18} & \num{31.88} & \num{6.90} & \num{32.39} \\
\quad \% Population     & & \num{41.77} & \num{5.61} & \num{21.56} & \num{28.94} & \num{0.73} & \num{1.38} \\
\quad Total Municipalities & & \num{215} & \num{728} & \num{1391} & \num{2581} & \num{559} & \num{2623} \\
\quad Total Population  & & \num{19735159} & \num{2649282} & \num{10187357} & \num{13674622} & \num{345513} & \num{652442} \\
\bottomrule
\end{tabular}
\end{adjustbox}
\noindent\footnotesize\emph{Notes.} Percentages are \emph{column shares}: \(100\%\) equals each cluster.
NLBP totals (CL1--CL3): population \num{32571798} (\pnum{68.94322975}\%), municipalities \num{2334} (\pnum{28.82549092}\%).
LBP totals (CL4--CL6): population \num{14672577} (\pnum{31.05677025}\%), municipalities \num{5763} (\pnum{71.17450908}\%).
\end{table}


\clearpage
\pagebreak



\begin{table}%[p]
\centering
\caption{Distribution by cluster within regions (100\% = CCAA; population and municipalities)}
\label{tbl:6}
\footnotesize
\begin{adjustbox}{max totalsize={\textwidth}{\textheight},keepaspectratio}
\begin{tabular}{llrrrrrrrrr}
\toprule
\multicolumn{2}{c}{} &
\multicolumn{6}{c}{\textbf{Clusters}} &
\multicolumn{2}{c}{\textbf{Aggregates}} &
\multicolumn{1}{c}{\textbf{Total}}\\
\cmidrule(lr){3-8}\cmidrule(lr){9-10}\cmidrule(lr){11-11}
\textbf{CCAA} &  &
\textbf{CL1} & \textbf{CL2} & \textbf{CL3} & \textbf{CL4} & \textbf{CL5} & \textbf{CL6} &
\textbf{NLBP} & \textbf{LBP} &
\textbf{(mun / pop)} \\
\midrule

\multirow{2}{*}{\textbf{Andalusia}} & pct\_mun (\%) & \pnum{3.25} & \pnum{2.08} & \pnum{5.46} & \pnum{84.92} & \pnum{2.34} & \pnum{1.95} & \pnum{10.79} & \pnum{89.21} & \num{769} \\
 & pct\_pop (\%) & \pnum{28.52} & \pnum{1.11} & \pnum{10.13} & \pnum{60.09} & \pnum{0.10} & \pnum{0.05} & \pnum{39.76} & \pnum{60.24} & \num{8474223} \\
\midrule

\multirow{2}{*}{\textbf{Aragon}} & pct\_mun (\%) & \pnum{0.14} & \pnum{7.53} & \pnum{6.99} & \pnum{17.67} & \pnum{24.93} & \pnum{42.74} & \pnum{14.66} & \pnum{85.34} & \num{730} \\
 & pct\_pop (\%) & \pnum{0.02} & \pnum{13.17} & \pnum{63.45} & \pnum{13.05} & \pnum{7.70} & \pnum{2.61} & \pnum{76.64} & \pnum{23.36} & \num{1323488} \\
\midrule

\multirow{2}{*}{\textbf{Asturias}} & pct\_mun (\%) & \pnum{2.56} & \pnum{11.54} & \pnum{7.69} & \pnum{35.90} & \pnum{10.26} & \pnum{32.05} & \pnum{21.79} & \pnum{78.21} & \num{78} \\
 & pct\_pop (\%) & \pnum{34.20} & \pnum{4.45} & \pnum{30.77} & \pnum{23.92} & \pnum{0.42} & \pnum{6.25} & \pnum{69.42} & \pnum{30.59} & \num{1004686} \\
\midrule

\multirow{2}{*}{\textbf{Balearic Islands}} & pct\_mun (\%) & \pnum{2.99} & \pnum{0} & \pnum{49.25} & \pnum{43.28} & \pnum{1.49} & \pnum{2.99} & \pnum{52.24} & \pnum{47.76} & \num{67} \\
 & pct\_pop (\%) & \pnum{39.66} & \pnum{0} & \pnum{35.59} & \pnum{24.65} & \pnum{0.05} & \pnum{0.04} & \pnum{75.25} & \pnum{24.74} & \num{1176659} \\
\midrule

\multirow{2}{*}{\textbf{Canary Islands}} & pct\_mun (\%) & \pnum{4.60} & \pnum{0} & \pnum{18.39} & \pnum{74.71} & \pnum{0} & \pnum{2.30} & \pnum{22.99} & \pnum{77.01} & \num{87} \\
 & pct\_pop (\%) & \pnum{28.99} & \pnum{0} & \pnum{20.21} & \pnum{50.52} & \pnum{0} & \pnum{0.28} & \pnum{49.20} & \pnum{50.80} & \num{2175730} \\
\midrule

\multirow{2}{*}{\textbf{Cantabria}} & pct\_mun (\%) & \pnum{4.90} & \pnum{12.75} & \pnum{29.41} & \pnum{28.43} & \pnum{0.98} & \pnum{23.53} & \pnum{47.06} & \pnum{52.94} & \num{102} \\
 & pct\_pop (\%) & \pnum{44.10} & \pnum{15.94} & \pnum{30.31} & \pnum{7.09} & \pnum{0.23} & \pnum{2.33} & \pnum{90.35} & \pnum{9.65} & \num{585402} \\
\midrule

\multirow{2}{*}{\textbf{Castile--La Mancha}} & pct\_mun (\%) & \pnum{0.11} & \pnum{4.25} & \pnum{9.48} & \pnum{34.86} & \pnum{13.29} & \pnum{38.02} & \pnum{13.84} & \pnum{86.17} & \num{918} \\
 & pct\_pop (\%) & \pnum{1.72} & \pnum{6.43} & \pnum{30.31} & \pnum{56.38} & \pnum{2.89} & \pnum{2.27} & \pnum{38.46} & \pnum{61.54} & \num{2050565} \\
\midrule

\multirow{2}{*}{\textbf{Castile and León}} & pct\_mun (\%) & \pnum{0.27} & \pnum{5.70} & \pnum{10.24} & \pnum{15.31} & \pnum{5.65} & \pnum{62.84} & \pnum{16.21} & \pnum{83.80} & \num{2247} \\
 & pct\_pop (\%) & \pnum{32.43} & \pnum{7.09} & \pnum{28.42} & \pnum{17.23} & \pnum{4.42} & \pnum{10.41} & \pnum{67.94} & \pnum{32.06} & \num{2369530} \\
\midrule

\multirow{2}{*}{\textbf{Catalonia}} & pct\_mun (\%) & \pnum{7.82} & \pnum{19.98} & \pnum{38.58} & \pnum{17.76} & \pnum{2.85} & \pnum{13.00} & \pnum{66.38} & \pnum{33.61} & \num{946} \\
 & pct\_pop (\%) & \pnum{67.11} & \pnum{8.15} & \pnum{19.90} & \pnum{4.00} & \pnum{0.21} & \pnum{0.62} & \pnum{95.16} & \pnum{4.83} & \num{7786738} \\
\midrule

\multirow{2}{*}{\textbf{Valencian Community}} & pct\_mun (\%) & \pnum{8.13} & \pnum{14.79} & \pnum{25.14} & \pnum{35.49} & \pnum{0.92} & \pnum{15.53} & \pnum{48.06} & \pnum{51.94} & \num{542} \\
 & pct\_pop (\%) & \pnum{40.99} & \pnum{11.66} & \pnum{19.35} & \pnum{27.35} & \pnum{0.06} & \pnum{0.58} & \pnum{72.00} & \pnum{27.99} & \num{5096865} \\
\midrule

\multirow{2}{*}{\textbf{Extremadura}} & pct\_mun (\%) & \pnum{0} & \pnum{1.83} & \pnum{0.52} & \pnum{80.10} & \pnum{2.88} & \pnum{14.66} & \pnum{2.35} & \pnum{97.64} & \num{382} \\
 & pct\_pop (\%) & \pnum{0} & \pnum{1.26} & \pnum{9.19} & \pnum{86.14} & \pnum{1.12} & \pnum{2.30} & \pnum{10.45} & \pnum{89.56} & \num{1046982} \\
\midrule

\multirow{2}{*}{\textbf{Galicia}} & pct\_mun (\%) & \pnum{0.64} & \pnum{9.65} & \pnum{8.04} & \pnum{55.95} & \pnum{2.25} & \pnum{23.47} & \pnum{18.33} & \pnum{81.67} & \num{311} \\
 & pct\_pop (\%) & \pnum{20.04} & \pnum{8.55} & \pnum{27.20} & \pnum{39.59} & \pnum{0.26} & \pnum{4.34} & \pnum{55.79} & \pnum{44.19} & \num{2679520} \\
\midrule

\multirow{2}{*}{\textbf{Community of Madrid}} & pct\_mun (\%) & \pnum{9.50} & \pnum{3.35} & \pnum{72.63} & \pnum{9.50} & \pnum{0} & \pnum{5.03} & \pnum{85.48} & \pnum{14.53} & \num{179} \\
 & pct\_pop (\%) & \pnum{78.73} & \pnum{0.52} & \pnum{19.89} & \pnum{0.84} & \pnum{0} & \pnum{0.02} & \pnum{99.14} & \pnum{0.86} & \num{6750336} \\
\midrule

\multirow{2}{*}{\textbf{Region of Murcia}} & pct\_mun (\%) & \pnum{2.22} & \pnum{4.44} & \pnum{8.89} & \pnum{84.44} & \pnum{0} & \pnum{0} & \pnum{15.55} & \pnum{84.44} & \num{45} \\
 & pct\_pop (\%) & \pnum{2.78} & \pnum{1.94} & \pnum{7.06} & \pnum{88.21} & \pnum{0} & \pnum{0} & \pnum{11.78} & \pnum{88.21} & \num{1531878} \\
\midrule

\multirow{2}{*}{\textbf{Navarre}} & pct\_mun (\%) & \pnum{2.94} & \pnum{22.79} & \pnum{27.57} & \pnum{9.56} & \pnum{11.03} & \pnum{26.10} & \pnum{53.30} & \pnum{46.69} & \num{272} \\
 & pct\_pop (\%) & \pnum{44.80} & \pnum{19.02} & \pnum{23.80} & \pnum{8.67} & \pnum{2.13} & \pnum{1.58} & \pnum{87.62} & \pnum{12.38} & \num{664117} \\
\midrule

\multirow{2}{*}{\textbf{Basque Country}} & pct\_mun (\%) & \pnum{8.80} & \pnum{28.80} & \pnum{54.00} & \pnum{3.20} & \pnum{3.60} & \pnum{1.60} & \pnum{91.60} & \pnum{8.40} & \num{250} \\
 & pct\_pop (\%) & \pnum{52.48} & \pnum{10.53} & \pnum{36.51} & \pnum{0.19} & \pnum{0.25} & \pnum{0.04} & \pnum{99.52} & \pnum{0.48} & \num{2207769} \\
\midrule

\multirow{2}{*}{\textbf{La Rioja}} & pct\_mun (\%) & \pnum{0.58} & \pnum{11.56} & \pnum{13.87} & \pnum{31.79} & \pnum{6.36} & \pnum{35.84} & \pnum{26.01} & \pnum{73.99} & \num{173} \\
 & pct\_pop (\%) & \pnum{46.90} & \pnum{15.04} & \pnum{22.75} & \pnum{11.33} & \pnum{2.11} & \pnum{1.87} & \pnum{84.69} & \pnum{15.31} & \num{319887} \\
\midrule
\bottomrule
\end{tabular}
\end{adjustbox}\\
\noindent\footnotesize\emph{Notes}. Percentages are row shares: \(100\%\) equals each CCAA. \textbf{NLBP} = CL1--CL3; \textbf{LBP} = CL4--CL6. Counts use thousands separators.

\end{table}
\pagebreak

\sisetup{
  group-separator={,},
  round-mode=places,
  round-precision=2,
  table-format=2.2
}

\pagebreak
\clearpage



\begin{sidewaystable}%[htbp]
\centering
\caption{Distribution of municipalities and population by degree of urbanisation within each cluster}
\label{tbl:8}
\footnotesize
\begin{tabular}{lrrrrrr|r}
\toprule
 & Cluster 1 & Cluster 2 & Cluster 3 & Cluster 4 & Cluster 5 & Cluster 6 & \textbf{Total} \\
\midrule
\textbf{Urban} & & & & & & & \\
\quad \% Municipalities & 61.35 & 3.98 & 23.11 & 11.55 &  &  & 100 \\
\quad \% Population     & 71.77 & 0.83 & 13.28 & 14.12 &  &  & 100 \\
\quad Municipalities & 154 & 10 & 58 & 29 &  &  & 251 \\
\quad Population    & 18,531,323 & 214,371 & 3,429,178 & 3,646,232 &  &  & 25,821,104 \\
\midrule
\textbf{Intermediate} & & & & & & & \\
\quad \% Municipalities & 5.35 & 16.31 & 36.41 & 41.94 &  &  & 100 \\
\quad \% Population     & 7.92 & 11.01 & 36.39 & 44.68 &  &  & 100 \\
\quad Municipalities  & 58 & 177 & 395 & 455 &  &  & 1,085 \\
\quad Population     & 1,202,093 & 1,670,184 & 5,521,741 & 6,779,311 &  &  & 15,173,329 \\
\midrule
\textbf{Rural} & & & & & & & \\
\quad \% Municipalities & 0.04 & 8.00 & 13.87 & 31.02 & 8.27 & 38.80 & 100 \\
\quad \% Population     & 0.03 & 12.24 & 19.78 & 51.99 & 5.53 & 10.44 & 100 \\
\quad Municipalities  & 3 & 541 & 938 & 2,097 & 559 & 2,623 & 6,761 \\
\quad Population    & 1,743 & 764,727 & 1,236,438 & 3,249,079 & 345,513 & 652,442 & 6,249,942 \\
\midrule
\textbf{Total} & & & & & & & \\
\quad Municipalities  & 215 & 728 & 1,391 & 2,581 & 559 & 2,623 & 8,097 \\
\quad Population    & 19,735,159 & 2,649,282 & 10,187,357 & 13,674,622 & 345,513 & 652,442 & 47,244,375 \\
\bottomrule
\end{tabular}
\end{sidewaystable}


\pagebreak
\clearpage

\begin{sidewaystable}%[htbp]
\centering
\caption{Distribution of municipalities and population across clusters within each degree of urbanisation}
\label{tbl:9}
\footnotesize
\begin{tabular}{lrrrrrr|r}
\toprule
 & Cluster 1 & Cluster 2 & Cluster 3 & Cluster 4 & Cluster 5 & Cluster 6 & \textbf{Total} \\
\midrule
\textbf{Urban} & & & & & & & \\
\quad \% Municipalities & 71.63 & 1.37 & 4.17 & 1.12 &   &   &   \\
\quad \% Population     & 93.90 & 8.09 & 33.66 & 26.66 &   &   &   \\
\quad Municipalities    & 154 & 10 & 58 & 29 &   &   & 251 \\
\quad Population        & 18,531,323 & 214,371 & 3,429,178 & 3,646,232 &   &   & 25,821,104 \\
\midrule
\textbf{Intermediate} & & & & & & & \\
\quad \% Municipalities & 26.98 & 24.31 & 28.40 & 17.63 &   &   &   \\
\quad \% Population     & 6.09 & 63.04 & 54.20 & 49.58 &   &   &   \\
\quad Municipalities    & 58 & 177 & 395 & 455 &   &   & 1,085 \\
\quad Population        & 1,202,093 & 1,670,184 & 5,521,741 & 6,779,311 &   &   & 15,173,329 \\
\midrule
\textbf{Rural} & & & & & & & \\
\quad \% Municipalities & 1.40 & 74.31 & 67.43 & 81.25 & 100 & 100 &   \\
\quad \% Population     & 0.01 & 28.87 & 12.14 & 23.76 & 100 & 100 &   \\
\quad Municipalities    & 3 & 541 & 938 & 2,097 & 559 & 2,623 & 6,761 \\
\quad Population        & 1,743 & 764,727 & 1,236,438 & 3,249,079 & 345,513 & 652,442 & 6,249,942 \\
\midrule
\textbf{Total} & & & & & & & \\
\quad Municipalities    & 215 & 728 & 1,391 & 2,581 & 559 & 2,623 & 8,097 \\
\quad Population        & 19,735,159 & 2,649,282 & 10,187,357 & 13,674,622 & 345,513 & 652,442 & 47,244,375 \\
\bottomrule
\end{tabular}
\end{sidewaystable}



\clearpage
\pagebreak


% \begin{table}%[htbp]
% \centering
% \caption{ANOVA Results by Variable, Grouped by Type}
% \label{tbl:10}
% \footnotesize
% \begin{tabular}{lrr}
% \toprule
% \textbf{Variable} & \textbf{$F$-value} & \textbf{$p$-value} \\
% \midrule
% \multicolumn{3}{l}{\textit{Socio-demographic variables}} \\
% \midrule
% Elderly dependency ratio          & 3,212.0  &  0.0000  \\
% Youth out-migration share         & 1,412.0  &  0.0000  \\
% Net migration rate                & 3.14    & 0.0763     \\
% Regional population growth        & 529.10   & 0.0000 \\
% Population density                & 10.37   & 0.0013    \\
% \midrule
% \multicolumn{3}{l}{\textit{Economic and financial variables}} \\
% \midrule
% Industry employment share         & 679.4   & 0.0000  \\
% Agricultural employment share    & 1,606.0  &  0.0000 \\
% Relative income per capita        & 1,431.0  & 0.0000 \\
% Bank branches per capita          & 5.19    & 0.0227    \\
% \midrule
% \multicolumn{3}{l}{\textit{Accessibility and connectivity variables}} \\
% \midrule
% Distance to urban and semiurban centers      & 2,457.0  & 0.0000 \\
% Market potential by income        & 523.0   & 0.0000 \\
% \bottomrule
% \end{tabular}
% % \begin{tablenotes}
% % \small
% % \item \textit{Significance codes:} \textbf{***} p $<$ 0.001; \textbf{**} p $<$ 0.01; \textbf{*} p $<$ 0.05; \textbf{.} p $<$ 0.1
% % \end{tablenotes}
% \end{table}

\begin{table}[htbp]
\centering
\caption{ANOVA Results by variable, grouped by type}
\label{tbl:anova}

\begin{threeparttable}
\footnotesize

\begin{tabular}{l S[table-format=4.1] S[table-format=1.3]}
\toprule
\textbf{Variable} & {\textbf{$F$-value}} & {\textbf{$p$-value}} \\
\midrule
\multicolumn{3}{l}{\textit{Socio-demographic variables}} \\
\midrule
Elderly dependency ratio          & 1257.0 & {<0.001} \\
Population under 16               & 414.4  & {<0.001} \\
Net migration rate                & 25.6   & {<0.001} \\
Regional population growth        & 1063.0 & {<0.001} \\
Population density                & 7008.0 & {<0.001} \\
\midrule
\multicolumn{3}{l}{\textit{Economic and financial variables}} \\
\midrule
Industrial employment share       & 3707.0 & {<0.001} \\
Agricultural employment share     & 693.6  & {<0.001} \\
Relative income per capita        & 828.9  & {<0.001} \\
Bank branches per capita          & 3723.0 & {<0.001} \\
\midrule
\multicolumn{3}{l}{\textit{Accessibility and connectivity variables}} \\
\midrule
Distance to urban centres         & 1180.0 & {<0.001} \\
Income market potential           & 645.7  & {<0.001} \\
\bottomrule
\end{tabular}

\begin{tablenotes}
\footnotesize
\item Note: All variables are standardised. p-values below 0.001 are reported as <0.001.
\end{tablenotes}

\end{threeparttable}
\end{table}
\clearpage
\pagebreak


%\begin{table}[H]
%\centering
%\caption{Average Marginal Effects (AME) from the Logistic Regression Model, by Variable Type}
%\label{tbl:9}
%\begin{adjustbox}{max width=\textwidth}
%\begin{tabular}{lrrrrrr}
%\toprule
%\textbf{Variable} & \textbf{AME} & \textbf{SE} & \textbf{z} & \textbf{p-value} & \textbf{95\% CI Lower} & \textbf{95\% CI Upper} \\
%\midrule
%\multicolumn{7}{l}{\textit{Socio-demographic variables}} \\
%\cmidrule(lr){1-7}
%Population density (inhabitants/km\(^2\)) & -0.0001 & 0.0000 & -7.71 & 0.0000 & -0.0001 & -0.0000 \\
%Elderly dependency ratio                  & 0.2056 & 0.0147 & 13.99 & 0.0000 & 0.1768 & 0.2344 \\
%Youth out-migration share                  & -0.2816 & 0.0205 & -13.76 & 0.0000 & -0.3217 & -0.2415 \\
%Regional population growth (2000–2022)     & -0.1951 & 0.0075 & -26.11 & 0.0000 & -0.2097 & -0.1804 \\
%Net migration rate                         & -0.4907 & 0.0614 & -7.99 & 0.0000 & -0.6111 & -0.3703 \\
%\midrule
%\multicolumn{7}{l}{\textit{Economic and financial variables}} \\
%\cmidrule(lr){1-7}
%Bank branches per capita                  & 0.0498 & 0.0043 & 11.47 & 0.0000 & 0.0413 & 0.0583 \\
%Agriculture employment share               & 0.2098 & 0.0121 & 17.36 & 0.0000 & 0.1861 & 0.2335 \\
%Industry employment share                  & -0.5988 & 0.0124 & -48.13 & 0.0000 & -0.6232 & -0.5744 \\
%Relative income per capita                  & -0.3641 & 0.0150 & -24.22 & 0.0000 & -0.3936 & -0.3347 \\
%\midrule
%\multicolumn{7}{l}{\textit{Accessibility and connectivity variables}} \\
%\cmidrule(lr){1-7}
%Distance to urban and semiurban centers      & 0.0051 & 0.0003 & 14.76 & 0.0000 & 0.0044 & 0.0058 \\
%Market potential by income                 & -0.0001 & 0.0000 & -25.83 & 0.0000 & -0.0001 & -0.0001 \\
%\bottomrule
%\end{tabular}
%\end{adjustbox}
%\end{table}




\clearpage
\pagebreak

\begin{table}[!ht]
\centering
\caption{Closest decision boundary by dimension}
\label{tab:threshold_dimension}
\footnotesize
\begin{tabular}{lrrrr}
\toprule
 & \multicolumn{2}{c}{\textbf{All LBP}} & \multicolumn{2}{c}{\textbf{Switch to LBP}} \\
\cmidrule(lr){2-3} \cmidrule(lr){4-5}
\textbf{Dimension} & \textbf{Freq.} & \textbf{\%} & \textbf{Freq.} & \textbf{\%} \\
\midrule
Geographic & 0    & 0.00  & 0   & 0.00  \\
Demographic   & 2,514 & 43.67 & 228 & 21.33 \\
Socioeconomic & 3,243 & 56.33 & 841 & 78.67 \\
\bottomrule
\end{tabular}

\vspace{0.2cm}
\begin{minipage}{0.9\linewidth}
\footnotesize
\textit{Note:} The table reports the dimension associated with the closest decision boundary for municipalities classified as left-behind (LBP) and for those that switch to left-behind status when geographical and accessibility-related variables are introduced. Percentages are computed within each group. Accessibility-related variables never define the closest decision boundary.
\end{minipage}
\end{table}

\begin{table}[!ht]
\centering
\caption{Structural distance to decision boundaries by type of left-behind place}
\label{tab:structural_blocks}
\footnotesize
\begin{tabular}{lccc}
\toprule
\textbf{Dimension} & \textbf{Cluster 4} & \textbf{Cluster 5} & \textbf{Cluster 6} \\
\midrule
Demographic    & 12.748 & 22.233 & 5.341 \\
Socioeconomic  & 7.706  & 10.771 & 20.637 \\
Geographic  & 5.167  & 7.068  & 5.911 \\
\bottomrule
\end{tabular}

\vspace{0.2cm}
\begin{minipage}{0.9\linewidth}
\footnotesize
\textit{Note:} Entries report the median, within each left-behind cluster, of municipality-level median absolute distance-to-threshold values (\(|dthr|\)) computed across variables in each dimension. For left-behind municipalities, thresholds are defined with respect to the closest non-left-behind cluster. Larger values indicate dimensions along which municipalities are structurally farther from any marginal transition out of left-behind status.

\end{minipage}
\end{table}
\begin{table}[ht]
\centering
\caption{Median distance to threshold by variable and LBP cluster}
\label{tab:threshold_by_cluster}
\begin{tabular}{lccc}
\toprule
\textbf{Variable} & \textbf{Cluster 4} & \textbf{Cluster 5} & \textbf{Cluster 6} \\
\midrule
Relative population growth        & 3.017 & 5.948 & 2.689 \\
Net migration rate                & 13.195 & 43.433 & 14.466 \\
Share of foreign-born under 16    & 51.087 & 98.746 & 5.293 \\
Elderly dependency ratio          & 9.110 & 9.552 & 2.852 \\
Population density                & 13.136 & 22.233 & 10.847 \\
Relative income per capita        & 2.621 & 12.059 & 15.545 \\
Industry employment share         & 33.991 & 322.366 & 11.834 \\
Agriculture employment share      & 3.267 & 13.721 & 119.583 \\
Bank branches per capita          & 14.680 & 2.604 & 24.562 \\
Distance to urban centres         & 7.507 & 5.558 & 3.116 \\
Market potential (PM\_renda)      & 2.830 & 8.592 & 8.704 \\
\bottomrule
\end{tabular}
\end{table}
\begin{table}[ht]
\centering
\footnotesize
\caption{Variables ordered by proximity to the classification frontier (median $|dthr|$)}
\label{tab:frontier_order}
\begin{tabular}{c p{4.2cm} p{4.2cm} p{4.2cm}}
\toprule
\textbf{Rank} & \textbf{Cluster 4} & \textbf{Cluster 5} & \textbf{Cluster 6} \\
\midrule
1  & Relative income (2.621)              & Bank branches (2.604)             & Population growth (2.689) \\
2  & Market potential (2.830)             & Distance to centres (5.558)       & Elderly dependency (2.852) \\
3  & Population growth (3.017)            & Population growth (5.948)         & Distance to centres (3.116) \\
4  & Agriculture share (3.267)            & Market potential (8.592)          & Foreign-born under 16 (5.293) \\

\bottomrule
\end{tabular}
\end{table}

\begin{sidewaystable}%[!ht]
    \centering
    \caption{Multinomial logit estimates and odds ratios (base: NLBP clusters)}
    \label{tbl:11}
    \footnotesize
    \begin{adjustbox}{max width=\textwidth}
    \begin{tabular}{p{5.5cm}*{6}{>{\raggedleft\arraybackslash}p{1.8cm}}}
        \toprule
        \multirow{2}{*}{\textbf{Variable}} & \multicolumn{2}{c}{\textbf{Category 1 (Cluster 4)}} & \multicolumn{2}{c}{\textbf{Category 2 (Cluster 5)}} & \multicolumn{2}{c}{\textbf{Category 3 (Cluster 6)}} \\
        \cmidrule(r){2-3}
        \cmidrule(r){4-5}
        \cmidrule(r){6-7}
        & Estimate & Odds Ratio & Estimate & Odds Ratio & Estimate & Odds Ratio \\
        \midrule
        \multicolumn{7}{l}{\textit{Socio-demographic}} \\
        \midrule
        Elderly dependency ratio           & 0.430* & 1.537 & 2.379*** & 10.789 & 5.895*** & 363.272 \\
        Share of foreign-born under 16     & 0.731*** & 2.078 & 0.318* & 1.374 & -3.520*** & 0.030 \\
        Net migration rate                 & -1.140*** & 0.320 & -0.715*** & 0.489 & -1.380*** & 0.252 \\
        Relative population growth & -5.360*** & 0.005 & -5.738*** & 0.003 & -7.035*** & 0.001 \\
        Population density (inhab/km\(^2\)) & -2.592*** & 0.075 & -11.341*** & 0.000 & -0.389 & 0.678 \\
        % \addlinespace
        \midrule
        \multicolumn{7}{l}{\textit{Economic and Financial}} \\
        \midrule
        Relative income per capita         & -6.127*** & 0.002 & -3.415*** & 0.033 & -0.910** & 0.403 \\
        Industry employment share          & -5.192*** & 0.006 & -6.640*** & 0.001 & -6.972*** & 0.001 \\
        Agriculture employment share       & 5.002*** & 148.662 & 2.741*** & 15.494 & -0.171 & 0.843 \\
        Bank branches per capita           & 1.069*** & 2.913 & 13.197*** & 538694 & -0.944* & 0.389 \\
        % \addlinespace
        \midrule
        \multicolumn{7}{l}{\textit{Second-Nature-Geography}} \\
        \midrule
        Distance to urban/semiurban centers & 0.649*** & 1.913 & 5.244*** & 189.500 & 5.268*** & 194.102 \\
        Market potential (PM\_renda)       & -5.864*** & 0.003 & -3.637*** & 0.026 & -2.353*** & 0.095 \\
        %\addlinespace
        \midrule
        Intercept                  & 1.258*** & 3.518 & -17.132*** & 0.000 & -0.274 & 0.760 \\
        \bottomrule
        \addlinespace
        \multicolumn{7}{p{17.3cm}}{\footnotesize\textit{Note:} Standardized estimates and odds ratios from a Multinomial Logit Model. The three clusters identified as Left-Behind Places (Clusters 4, 5, and 6) were recoded as categories 1, 2, and 3. The remaining clusters (1, 2, and 3), corresponding to municipalities not classified as LBPs, were grouped together and coded as 0, serving as the reference category. Significance levels: *** \(p<0.01\), ** \(p<0.05\), * \(p<0.1\). Source: Authors' own elaboration.} \\
    \end{tabular}
    \end{adjustbox}
\end{sidewaystable}




\pagebreak
\clearpage






\begin{figure}%[ht]
    \centering
    \caption{Spain's municipalities grouped by Autonomous Community.}
    \label{fig:map_ccaa_municipal_labels}\includegraphics[width=0.85\textwidth]{Figures/mapa_ccaa.png}

\end{figure}



\pagebreak
\clearpage


\begin{figure}%[ht]
    \centering
     \caption{Spain municipalities by level of urbanization}
    \label{fig:urbanizacion_municipios}\includegraphics[width=0.85\textwidth]{Figures/degurba.png}

\end{figure}


\pagebreak
\clearpage



\begin{figure}
  \centering
\caption{\centering Cluster classification of Spanish municipalities in 2022 (excluding geographical and distance variables)}
\footnotesize{Silhouette coefficient: 0.34}
  \label{fig1:clustermap}
  \begin{center}
\includegraphics[width=0.75\linewidth]{Figures/mapa_34_Dens_NoDist.png}
  \end{center}
\end{figure}


\clearpage
\pagebreak


\begin{figure}
  \centering
\caption{\centering Cluster classification of Spanish municipalities in 2022 (including geographical and distance variables)}
\footnotesize{Silhouette coefficient: 0.32}
  \label{fig2:clustermap}
  \begin{center}
\includegraphics[width=0.75\linewidth]{Figures/mapa32.png}
  \end{center}
\end{figure}


\clearpage
\pagebreak



\begin{figure}
  \centering
\caption{\centering Municipalities changing clusters with the inclusion of geographical and distance variables, 2022}
\footnotesize{\footnotesize{\centering Municipalities reclassified from clusters without geographical and distance variables to clusters with geographical and distance variables.}}
  \label{fig3:clustermap}
  \begin{center}
\includegraphics[width=0.75\linewidth]{Figures/dif32.png}
\end{center}
\end{figure}

\clearpage
\pagebreak

\begin{figure}[t]
\centering
\caption{Relationship between the share of municipalities classified as LBPs and the share of population residing in LBP municipalities across Spanish regions.}
\label{fig:lbp_scatter}
\includegraphics[width=0.8\textwidth]{Figures/dispersion.png}

\end{figure}

\clearpage
\pagebreak

%\begin{figure}[!ht]
%\centering
%\includegraphics[width=0.75\textwidth]{Figures/box.png}
%\caption{Distribution of absolute distance to the closest decision boundary by dimension}
%\label{fig:dthr_dimension}
%\footnotesize
%\textit{Note:} The figure displays the distribution of the absolute distance to the closest decision boundary (\(|dthr|\)) for municipalities classified as left-behind, grouped by the dimension that defines their nearest cluster boundary. Smaller values indicate closer proximity to a decision boundary. Accessibility-related variables do not appear because they never constitute the closest decision boundary for left-behind municipalities.
%\end{figure}




\begin{sidewaysfigure}
  \centering
\caption{\centering Estimated odds ratios from the multinomial logit model by type of left-behind place}
  \label{fig4:odds}
  \begin{center}
\includegraphics[width=.85\linewidth]{Figures/odds_tres.png}
\end{center}
\end{sidewaysfigure}


\clearpage
\pagebreak










\end{document}
